\documentclass[11pt,a4paper,oneside,openany]{book}

\usepackage[T1]{fontenc}
\usepackage[utf8]{inputenc}
\usepackage{libertinus}
\usepackage{microtype}
\usepackage[a4paper,margin=27mm,headheight=15pt]{geometry}
\usepackage{amsmath,amssymb,amsthm,mathtools}
\usepackage{bm,mathrsfs}
\usepackage{booktabs,longtable,array,tabularx,multirow}
\usepackage{enumitem}
\usepackage{xcolor}
\usepackage{tikz}
\usetikzlibrary{arrows.meta,positioning,calc,shapes.geometric}
\usepackage[most]{tcolorbox}
\usepackage{fancyhdr}
\usepackage{xurl}
\usepackage{listings}
\usepackage{caption}
\usepackage{hyperref}
\usepackage[nameinlink,noabbrev,capitalise]{cleveref}

\definecolor{deepolive}{RGB}{67,79,45}
\definecolor{softolive}{RGB}{236,239,226}
\definecolor{darkteal}{RGB}{31,83,84}
\definecolor{warmgray}{RGB}{92,87,78}
\definecolor{softcoral}{RGB}{246,225,219}
\definecolor{linkblue}{RGB}{31,76,122}
\definecolor{codeback}{RGB}{247,247,243}

\hypersetup{
  colorlinks=true,
  linkcolor=deepolive,
  citecolor=darkteal,
  urlcolor=linkblue,
  pdftitle={The Butcher-Tree Counting Transseries and Its Asymptotic Inverse},
  pdfauthor={A self-contained derivation prepared for Vladimir Reshetnikov},
  pdfsubject={A000081, Polya trees, rooted trees, Lambert W, Puiseux expansions, inverse transseries}
}

\setlist{itemsep=2pt,topsep=4pt,parsep=1pt}
\allowdisplaybreaks[3]
\raggedbottom
\emergencystretch=2em
\setcounter{tocdepth}{2}
\setcounter{secnumdepth}{3}

\pagestyle{fancy}
\fancyhf{}
\fancyhead[L]{\small\nouppercase{\leftmark}}
\fancyhead[R]{\small\thepage}
\renewcommand{\headrulewidth}{0.4pt}
\fancypagestyle{plain}{%
  \fancyhf{}%
  \fancyfoot[C]{\thepage}%
  \renewcommand{\headrulewidth}{0pt}%
}

\newtheorem{theorem}{Theorem}[chapter]
\newtheorem{lemma}[theorem]{Lemma}
\newtheorem{proposition}[theorem]{Proposition}
\newtheorem{corollary}[theorem]{Corollary}
\newtheorem{identity}[theorem]{Identity}
\newtheorem{conjecture}[theorem]{Conjecture}
\theoremstyle{definition}
\newtheorem{definition}[theorem]{Definition}
\newtheorem{example}[theorem]{Example}
\newtheorem{algorithm}[theorem]{Algorithm}
\theoremstyle{remark}
\newtheorem{remark}[theorem]{Remark}
\newtheorem{historical}[theorem]{Historical note}

\newtcolorbox{masterbox}[1][]{
  enhanced,breakable,colback=softolive,colframe=deepolive,
  boxrule=0.7pt,arc=1mm,left=2mm,right=2mm,top=1.5mm,bottom=1.5mm,
  title={Master formula},fonttitle=\bfseries,#1
}
\newtcolorbox{warningbox}[1][]{
  enhanced,breakable,colback=softcoral,colframe=warmgray,
  boxrule=0.6pt,arc=1mm,left=2mm,right=2mm,top=1.5mm,bottom=1.5mm,
  title={Important qualification},fonttitle=\bfseries,#1
}
\newtcolorbox{resultbox}[1][]{
  enhanced,breakable,colback=white,colframe=darkteal,
  boxrule=0.8pt,arc=1mm,left=2mm,right=2mm,top=1.5mm,bottom=1.5mm,
  title={Principal result},fonttitle=\bfseries,#1
}

\lstset{
  basicstyle=\ttfamily\small,
  backgroundcolor=\color{codeback},
  frame=single,
  rulecolor=\color{warmgray},
  breaklines=true,
  columns=fullflexible,
  keepspaces=true,
  showstringspaces=false,
  keywordstyle=\bfseries\color{darkteal},
  commentstyle=\itshape\color{warmgray}
}

\newcommand{\N}{\mathbb N}
\newcommand{\C}{\mathbb C}
\newcommand{\R}{\mathbb R}
\newcommand{\e}{\mathrm e}
\newcommand{\ii}{\mathrm i}
\newcommand{\dd}{\,\mathrm d}
\newcommand{\Coeff}[3]{[ #1^{#2}]\,#3}
\newcommand{\Res}{\operatorname*{Res}}
\newcommand{\Bell}[1]{\mathcal B_{#1}}
\newcommand{\pBell}[2]{\mathcal B_{#1,#2}}
\newcommand{\falling}[2]{#1^{\underline{#2}}}
\newcommand{\bigO}{\mathcal O}
\newcommand{\asympseries}{\sim_{\!\mathrm{asy}}}
\newcommand{\LambertW}{W}
\newcommand{\Tree}{\mathcal T}
\newcommand{\Action}{\mathfrak A}
\newcommand{\Stokes}{\mathfrak S}
\DeclareMathOperator{\Log}{Log}
\DeclareMathOperator{\sgn}{sgn}

\title{\Huge\bfseries The Butcher-Tree Counting Transseries\\[3mm]
\Large All-Order P\'olya-Tree Asymptotics,\\
Lambert--$W$ Reversion, and the Asymptotic Inverse of A000081}
\author{A self-contained derivation prepared for Vladimir Reshetnikov}
\date{September 2026}

\begin{document}
\frontmatter
\hypersetup{pageanchor=false}
\maketitle
\hypersetup{pageanchor=true}

\chapter*{Abstract}
\addcontentsline{toc}{chapter}{Abstract}
Let $a_n$ be OEIS A000081, the number of unlabelled, non-plane rooted trees with
$n$ vertices.  These are the rooted-tree shapes that index classical Butcher
series and B-series.  Their ordinary generating function
\[
 T(z)=\sum_{n\ge1}a_nz^n
\]
satisfies P\'olya's implicit multiset equation
\[
 T(z)=z\exp\!\left(\sum_{k\ge1}\frac{T(z^k)}{k}\right).
\]
This article develops the complete coefficient calculus associated with that
equation.  We factor the generating function through the Cayley tree function,
\[
 T(z)=-W_0\!\left(-z\exp\!\left(\sum_{k\ge2}\frac{T(z^k)}{k}\right)\right),
\]
and exploit the universal square-root singularity of $W_0$ at $-1/e$.

The first principal result is an all-orders Puiseux expansion at the dominant
singularity $\rho$.  Every coefficient is given explicitly by a finite formula
in universal Lambert coefficients, complete and partial Bell polynomials, and
the derivatives of the analytic P\'olya disturbance.  Exact singularity
transfer then gives
\[
 a_n\sim \rho^{-n}n^{-3/2}\sum_{k\ge0}\frac{A_k}{n^k},
\]
where every $A_k$ has a closed finite formula involving Bernoulli polynomials
and Bell polynomials.  Numerically,
\begin{align*}
 \rho&=0.33832185689920769519611262571701705318\ldots,\\
 \rho^{-1}&=2.95576528565199497471481752412319\ldots .
\end{align*}
and
\begin{align*}
 a_n\sim{}&0.439924012571025304040903391434545\ldots\,
       \rho^{-n}n^{-3/2}\\
 &\times\left(1+\frac{0.10040348009640143476\ldots}{n}+\cdots\right).
\end{align*}

The second principal result is the full asymptotic inverse.  After normalizing
$x$ by $X=(-\log\rho)x$, the inverse problem reduces to
\[
 Z=X-\frac32\log X+K(X^{-1}),
\]
where the coefficients of $K$ are explicit Bell-polynomial transforms of the
forward coefficients.  The entire logarithmic skeleton is summed exactly by
the real branch $W_{-1}$, and the remaining correction is an ordinary inverse
power series with a Lagrange--B\"urmann coefficient formula.  An equivalent
polynomial--logarithmic expansion is supplied with a closed recurrence for all
coefficient polynomials.

Finally, we distinguish the smooth asymptotic inverse from the integer
staircase inverse, derive the unavoidable periodic rounding layer, and give a
contour-resolved transseries over the recursively generated singularity atlas.
The first inherited singularity at $-\sqrt\rho$ is worked out explicitly.  It
produces an alternating forward sector of size $\rho^{-n/2}n^{-3/2}$ and, after
inversion, a log-periodic sector of magnitude
$y^{-1/2}(\log y)^{-3/4}$.  Thus both the perturbative logarithmic expansion and
the beyond-all-orders exponential sectors are described by general coefficient
formulae.

\chapter*{How to read the word ``full''}
\addcontentsline{toc}{chapter}{How to read the word ``full''}
There are three nested senses in which an asymptotic description can be full.
\begin{enumerate}
\item The \emph{dominant Poincar\'e sector} contains every inverse power of $n$
      multiplying $\rho^{-n}n^{-3/2}$.  In this article its coefficients are
      explicit and unconditional.
\item The \emph{smooth inverse sector} contains every polynomial in
      $\log\log y$ divided by a power of $\log y$, or equivalently every inverse
      power correction to an exact $W_{-1}$ core.  Its coefficients are also
      explicit and unconditional as formal asymptotic coefficients.
\item A \emph{global transseries} additionally resolves exponentially smaller
      contributions from other singularities of analytic continuation.  Since
      the P\'olya-tree generating function has many sheets and the unit circle
      is a natural boundary, such sectors must be indexed by a chosen
      continuation sheet and contour.  We give the complete recursive local
      calculus and the corresponding contour-resolved coefficient formula;
      Stokes/intersection multipliers are retained whenever they depend on that
      global choice.
\end{enumerate}
This distinction prevents the common but serious error of calling a single
finite numerical expansion a transseries, and also prevents a sheet-dependent
analytic-continuation statement from being presented as globally canonical.

\tableofcontents

\mainmatter

\part{Combinatorics and Coefficient Calculus}

\chapter{Butcher trees and the sequence A000081}
\label{ch:butcher}

\section{The objects being counted}
A \emph{rooted tree} is a finite connected acyclic graph with one distinguished
vertex.  It is \emph{non-plane} when the children of a vertex carry no left-to-right
order, and it is \emph{unlabelled} when two trees differing only by a relabelling of
vertices are identified.  Let $a_n$ be the number of isomorphism classes of such
trees on $n$ vertices.  We use
\[
 a_0=0,
 \qquad
 (a_n)_{n\ge1}=1,1,2,4,9,20,48,115,286,719,1842,\ldots .
\]
This is OEIS A000081 \cite{OEISA000081,Otter1948}.  The same tree shapes occur as the combinatorial indices of
Butcher series \cite{Butcher1963,Butcher1972,Butcher2021,McLachlan2017}: removing the root decomposes a tree into an unordered multiset of
rooted trees, while grafting the roots of a multiset of trees to a new vertex
reconstructs the original tree.

The adjective ``Butcher'' describes the role of these shapes in numerical analysis;
it does not change the enumeration.  A B-series assigns a coefficient and an
elementary differential to each non-plane rooted tree.  Consequently $a_n$ is the
number of possible tree shapes of order $n$ before method-dependent weights,
symmetry factors, and elementary differentials are attached.

\section{The multiset specification}
Let $\Tree$ denote the combinatorial class of unlabelled non-plane rooted trees and
let $\mathcal Z$ denote one atom.  The root decomposition is
\begin{equation}
 \Tree=\mathcal Z\times\operatorname{MSET}(\Tree).
 \label{eq:species}
\end{equation}
For unlabelled objects, the multiset construction is encoded by the cycle index of
the symmetric groups \cite{Polya1937,BergeronLabelleLeroux1998}.  Therefore the ordinary generating function
\begin{equation}
 T(z)=\sum_{n\ge1}a_nz^n
 \label{eq:Tdef}
\end{equation}
satisfies
\begin{equation}
 \boxed{
 T(z)=z\exp\!\left(\sum_{k\ge1}\frac{T(z^k)}{k}\right).}
 \label{eq:polya-functional}
\end{equation}
The identity holds formally in $\mathbb Z[[z]]$; no analytic convergence is needed
for its coefficientwise interpretation.

\begin{proof}[Direct derivation]
A multiset containing $m_j$ copies of each tree shape $j$ contributes the monomial
$z^{\sum_jm_j|j|}$.  For one available object of size $r$, arbitrary multiplicity
contributes $(1-z^r)^{-1}$.  Since there are $a_r$ possible shapes of size $r$, the
forest attached to the root has generating function
\[
 \prod_{r\ge1}(1-z^r)^{-a_r}.
\]
Multiplication by the root gives
\begin{equation}
 T(z)=z\prod_{r\ge1}(1-z^r)^{-a_r}.
 \label{eq:euler-product}
\end{equation}
Taking a formal logarithm and using
$-\log(1-z^r)=\sum_{k\ge1}z^{rk}/k$ yields
\[
 \log\frac{T(z)}z
 =\sum_{r\ge1}a_r\sum_{k\ge1}\frac{z^{rk}}k
 =\sum_{k\ge1}\frac1k\sum_{r\ge1}a_r(z^k)^r
 =\sum_{k\ge1}\frac{T(z^k)}k,
\]
which is \cref{eq:polya-functional}.
\end{proof}

\section{An exact quadratic-time recurrence}
Define the divisor transform
\begin{equation}
 b_m=\sum_{d\mid m}d\,a_d.
 \label{eq:bdivisor}
\end{equation}
Then the sequence is generated exactly by the following recurrence.

\begin{theorem}[Otter--P\'olya recurrence]
\label{thm:exact-recurrence}
For $n\ge2$,
\begin{equation}
 \boxed{
 (n-1)a_n=\sum_{j=1}^{n-1}b_j a_{n-j},
 \qquad b_j=\sum_{d\mid j}d\,a_d,}
 \label{eq:exact-recurrence}
\end{equation}
with $a_1=1$.
\end{theorem}

\begin{proof}
Set $F(z)=T(z)/z=\sum_{n\ge0}a_{n+1}z^n$.  From
\cref{eq:euler-product},
\[
 \log F(z)=\sum_{m\ge1}\frac{b_m}{m}z^m.
\]
Logarithmic differentiation gives
\[
 F'(z)=F(z)\sum_{m\ge1}b_mz^{m-1}.
\]
The coefficient of $z^{n-1}$ is
\[
 n a_{n+1}=\sum_{j=1}^{n}b_j a_{n-j+1}.
\]
Replacing $n$ by $n-1$ proves \cref{eq:exact-recurrence}.
\end{proof}

\begin{remark}[Why the recurrence is genuinely triangular]
Although $b_n$ contains $a_n$, the formula for $a_n$ uses only $b_j$ with
$j<n$.  After computing $a_n$, one may add the contribution $n a_n$ to every
$b_{kn}$.  This sieve implementation takes $O(N^2)$ integer arithmetic operations
for $a_1,\ldots,a_N$, and it avoids expansion of large products.
\end{remark}

\chapter{Bell polynomials, powers, and reversion}
\label{ch:calculus}

\section{Exponential Bell polynomials}
For integers $n\ge k\ge0$, the exponential partial Bell polynomial is
\begin{equation}
 \pBell nk(x_1,x_2,\ldots)
 =n!\!\sum_{\substack{m_1,m_2,\ldots\ge0\\
              \sum_jm_j=k,\ \sum_jjm_j=n}}
   \prod_{j\ge1}\frac1{m_j!}\left(\frac{x_j}{j!}\right)^{m_j}.
 \label{eq:partial-bell}
\end{equation}
The complete polynomial is
$\Bell n=\sum_{k=0}^n\pBell nk$.  We shall repeatedly use the generating identities
\begin{align}
 \frac1{k!}\left(\sum_{j\ge1}x_j\frac{t^j}{j!}\right)^k
 &=\sum_{n\ge k}\pBell nk(x_1,x_2,\ldots)\frac{t^n}{n!},
 \label{eq:bell-partial-egf}\\
 \exp\!\left(\sum_{j\ge1}x_j\frac{t^j}{j!}\right)
 &=\sum_{n\ge0}\Bell n(x_1,\ldots,x_n)\frac{t^n}{n!}.
 \label{eq:bell-complete-egf}
\end{align}
These are exactly the coefficient packages required by Fa\`a di Bruno composition,
exponentiation, and logarithmic conversion \cite{Comtet1974,ReshetnikovCCC}.

\section{A unit-series power formula}
Let
\[
 U(X)=1+\sum_{j\ge1}u_jX^j.
\]
For arbitrary $\beta\in\C$, formal binomial expansion and
\cref{eq:bell-partial-egf} give
\begin{equation}
 \boxed{
 \Coeff XN{U(X)^\beta}
 =\frac1{N!}\sum_{k=0}^N\falling\beta k\,
   \pBell Nk(1!u_1,2!u_2,\ldots),}
 \label{eq:unit-power-bell}
\end{equation}
where $\falling\beta k=\beta(\beta-1)\cdots(\beta-k+1)$.
Every sum in this article involving a fractional power of a unit series can
therefore be reduced to a finite Bell-polynomial expression.

\section{Lagrange--B\"urmann inversion}
We record the form of Lagrange inversion used below
\cite{Comtet1974,FlajoletSedgewick2009,ReshetnikovCCC}.

\begin{theorem}[Lagrange--B\"urmann]
\label{thm:lagrange}
Let $q=w\phi(w)$ with $\phi(0)\ne0$, and let $w=w(q)$ be the inverse germ with
$w(0)=0$.  Then, for $m\ge1$,
\begin{equation}
 \Coeff qm{w(q)}=\frac1m\Coeff u{m-1}{\phi(u)^{-m}}.
 \label{eq:lagrange-basic}
\end{equation}
More generally, for a formal equation $w=t\Phi(w)$,
\begin{equation}
 w(t)=\sum_{m\ge1}\frac{t^m}{m!}
 \left.\frac{\dd^{m-1}}{\dd u^{m-1}}\Phi(u)^m\right|_{u=0}.
 \label{eq:lagrange-derivative}
\end{equation}
\end{theorem}

\begin{proof}
Formal residue substitution gives
\[
 \Coeff qm{w(q)}=\Res_{q=0}\frac{w(q)}{q^{m+1}}\dd q.
\]
Substitute $q=u\phi(u)$ and integrate the derivative term formally.  The result is
$m^{-1}\Coeff u{m-1}\phi(u)^{-m}$.  Applying the same identity to
$w=t\Phi(w)$ and Taylor-expanding the coefficient gives
\cref{eq:lagrange-derivative}.
\end{proof}

\section{Logarithms as Bell transforms}
If $F(t)=1+\sum_{j\ge1}f_jt^j$, then
\begin{equation}
 \Coeff tn{\log F(t)}
 =\frac1{n!}\sum_{k=1}^n(-1)^{k-1}(k-1)!
   \pBell nk(1!f_1,2!f_2,\ldots).
 \label{eq:log-bell}
\end{equation}
This identity will convert the forward asymptotic coefficients into the additive
phase coefficients needed for inversion.

\begin{masterbox}
The entire calculation uses only four coefficient operations:
\begin{enumerate}
\item exact divisor convolution for the coefficients $a_n$;
\item complete Bell polynomials for exponentials and derivative jets;
\item partial Bell polynomials for arbitrary powers and logarithms;
\item Lagrange--B\"urmann inversion for the Lambert branch and for the inverse
      counting law.
\end{enumerate}
No coefficient is defined by an unnamed remainder or an implicit numerical fit.
\end{masterbox}

\part{The Dominant Forward Transseries}

\chapter{Cayley--Lambert factorization and the dominant singularity}
\label{ch:factorization}

\section{Separating the self-interaction}
Write
\begin{equation}
 R(z)=\sum_{k\ge2}\frac{T(z^k)}k,
 \qquad
 \zeta(z)=z\e^{R(z)}.
 \label{eq:R-zeta}
\end{equation}
Then \cref{eq:polya-functional} becomes
\begin{equation}
 T(z)=\zeta(z)\e^{T(z)}.
 \label{eq:T-zeta}
\end{equation}
Let $C(u)$ be the Cayley tree function, characterized near $u=0$ by
$C=u\e^C$.  Since $C(u)=-W_0(-u)$ \cite{Corless1996,Genitrini2016},
\begin{equation}
 \boxed{T(z)=C(\zeta(z))=-W_0(-\zeta(z)).}
 \label{eq:Lambert-factorization}
\end{equation}
The nontrivial P\'olya symmetries are now confined to the analytic disturbance
$R$; the singular mechanism itself is the universal Lambert branch point.

\section{A useful coefficient representation of the disturbance}
From the exact coefficients of $T$,
\begin{align}
 R(z)
 &=\sum_{k\ge2}\sum_{m\ge1}\frac{a_m}{k}z^{km}
   =\sum_{r\ge2}h_rz^r,\label{eq:R-coeff}\\
 h_r
 &=\frac1r\sum_{\substack{d\mid r\\d<r}}d\,a_d
   =\frac{b_r-r a_r}{r}.
 \label{eq:h-coeff}
\end{align}
Unlike $a_r$, the coefficient $h_r$ depends only on tree numbers with indices at
most $r/2$.  This is analytically decisive: $R$ has radius of convergence
$\sqrt\rho$, whereas $T$ has radius $\rho$.

\section{Characterization of the Otter singularity}
Let $\rho$ be the radius of convergence of $T$.  Positivity and aperiodicity of
$(a_n)$ imply that the positive real point $\rho$ is the unique singularity of
$T$ on $|z|=\rho$.  Since $R$ is analytic for $|z|<\sqrt\rho$, the first
singularity in \cref{eq:Lambert-factorization} occurs when its argument reaches
the Cayley critical value $1/e$.  Hence
\begin{equation}
 T(\rho)=1,
 \qquad
 \zeta(\rho)=\frac1e,
 \qquad
 \log\rho+1+R(\rho)=0.
 \label{eq:rho-characterization}
\end{equation}
The last equation is especially convenient numerically because the coefficients
$h_r$ decay geometrically after multiplication by $\rho^r$.

\begin{theorem}[Dominant square-root singularity]
\label{thm:dominant-singularity}
There is a unique $\rho\in(0,1)$ satisfying
\cref{eq:rho-characterization}.  In a dented neighbourhood of $\rho$, $T$ has a
convergent Puiseux expansion in powers of $(1-z/\rho)^{1/2}$.  Its leading term is
\begin{equation}
 T(z)=1-\sqrt{2\ell_1}\left(1-\frac z\rho\right)^{1/2}
      +O\!\left(1-\frac z\rho\right),
 \qquad
 \ell_1=1+\rho R'(\rho)>0.
 \label{eq:leading-singular}
\end{equation}
\end{theorem}

\begin{proof}
The classical P\'olya--Otter argument gives a finite positive radius and the
critical equations above \cite{Polya1937,Otter1948,FlajoletSedgewick2009}.  Differentiating $\zeta=z\e^R$ gives
$e\rho\zeta'(\rho)=1+\rho R'(\rho)=\ell_1>0$.  The local inverse of
$u\mapsto u\e^{-u}$ at its nondegenerate maximum $u=1$ therefore has a
square-root branch, equivalently $W_0$ has its standard square-root branch at
$-1/e$.  Substitution of the nonzero linear term of $1-e\zeta(z)$ gives
\cref{eq:leading-singular}.  Aperiodicity excludes another singularity of equal
modulus.
\end{proof}

\section{High-precision constants}
Using \cref{eq:exact-recurrence,eq:h-coeff,eq:rho-characterization} gives
\begin{align}
 \rho={}&0.3383218568992076951961126257170170531837746075329677955723038\ldots,
 \label{eq:rho-numeric}\\
 \alpha:=\rho^{-1}={}&2.9557652856519949747148175241231945883754923046635965953504725\ldots,
 \label{eq:alpha-numeric}\\
 \lambda:=\log\alpha={}&1.0837575972446246604827119808961742790496621242026722349144271\ldots,
 \label{eq:lambda-numeric}\\
 R(\rho)={}&0.0837575972446246604827119808961742790496621242026722349144271\ldots.
 \label{eq:Rrho-numeric}
\end{align}
The equality $\lambda=1+R(\rho)$ is an immediate rearrangement of
\cref{eq:rho-characterization}.

\chapter{The universal Lambert branch coefficients}
\label{ch:universal-lambert}

\section{A canonical square-root coordinate}
Set
\begin{equation}
 w=1-T,
 \qquad
 q=\sqrt{2(1-e\zeta)}.
 \label{eq:wq-def}
\end{equation}
Because $e\zeta=T\e^{1-T}=(1-w)\e^w$,
\begin{align}
 q^2
 &=2\bigl(1-(1-w)\e^w\bigr)\notag\\
 &=w^2 h(w),
 \qquad
 h(w)=\sum_{j\ge0}\frac{2}{(j+2)j!}w^j
      =1+\frac23w+\frac14w^2+\frac1{15}w^3+\cdots .
 \label{eq:q-w-h}
\end{align}
We choose the branch for which $q/w\to1$.  Then $q=w\sqrt{h(w)}$ is an ordinary
invertible power series.

\begin{theorem}[Universal Lambert coefficients]
\label{thm:omega}
The inverse series has the form
\begin{equation}
 w(q)=\sum_{m\ge1}\omega_mq^m,
 \qquad
 \boxed{
 \omega_m=\frac1m\Coeff u{m-1}{h(u)^{-m/2}}.}
 \label{eq:omega-master}
\end{equation}
Consequently
\begin{equation}
 -W_0(-\zeta)=1-\sum_{m\ge1}\omega_m
        \bigl(2(1-e\zeta)\bigr)^{m/2}.
 \label{eq:W-universal-expansion}
\end{equation}
\end{theorem}

\begin{proof}
Apply \cref{thm:lagrange} to $q=w\phi(w)$ with
$\phi(w)=\sqrt{h(w)}$.  This gives
$\omega_m=m^{-1}[u^{m-1}]\phi(u)^{-m}$, which is
\cref{eq:omega-master}.  Since $T=1-w$, substitution proves
\cref{eq:W-universal-expansion}.
\end{proof}

The first universal coefficients are
\begin{align}
 w(q)={}&q-\frac13q^2+\frac{11}{72}q^3-\frac{43}{540}q^4
 +\frac{769}{17280}q^5-\frac{221}{8505}q^6\notag\\
 &+\frac{680863}{43545600}q^7-\frac{1963}{204120}q^8
 +\frac{226287557}{37623398400}q^9
 -\frac{5776369}{1515591000}q^{10}+\cdots .
 \label{eq:omega-first}
\end{align}
They are rational and independent of the tree class.  Every class satisfying
$T=\zeta\e^T$ with analytic $\zeta$ uses exactly the same sequence
$(\omega_m)$; only the expansion of $1-e\zeta(z)$ changes.

\chapter{Complete Puiseux expansion at the dominant singularity}
\label{ch:puiseux}

\section{Scaled logarithmic jets}
Put
\begin{equation}
 X=1-\frac z\rho,
 \qquad
 L(z)=\log\zeta(z)=\log z+R(z),
 \label{eq:X-L}
\end{equation}
and define the scaled jets
\begin{equation}
 \ell_j=\rho^jL^{(j)}(\rho)
 =(-1)^{j-1}(j-1)!+\rho^jR^{(j)}(\rho).
 \label{eq:ell-jets}
\end{equation}
The derivatives of $R$ are themselves explicit convergent sums:
\begin{equation}
 \rho^jR^{(j)}(\rho)
 =\sum_{m\ge j}h_m\falling mj\rho^m.
 \label{eq:R-derivative-sum}
\end{equation}
Define
\begin{equation}
 S(X)=2\bigl(1-e\zeta(\rho(1-X))\bigr)
     =\sum_{j\ge1}s_jX^j.
 \label{eq:S-def}
\end{equation}
Since $e\zeta(\rho)=1$, complete Bell polynomials exponentiate the Taylor
series of $L+1$.

\begin{lemma}[Disturbance coefficients]
\label{lem:sj}
For $j\ge1$,
\begin{equation}
 \boxed{
 s_j=\frac{2(-1)^{j+1}}{j!}
 \Bell j(\ell_1,\ell_2,\ldots,\ell_j).}
 \label{eq:sj-bell}
\end{equation}
In particular $s_1=2\ell_1=2(1+\rho R'(\rho))>0$.
\end{lemma}

\begin{proof}
Taylor expansion at $z=\rho$ gives
\[
 1+L(\rho(1-X))
 =\sum_{j\ge1}(-1)^j\ell_j\frac{X^j}{j!}.
\]
Use \cref{eq:bell-complete-egf}.  Every monomial of
$\Bell j$ has weighted degree $j$, so replacing $\ell_r$ by
$(-1)^r\ell_r$ multiplies the polynomial by $(-1)^j$.  Since
$S=2(1-\exp(1+L))$, the stated sign follows.
\end{proof}

\section{The all-order coefficient formula}
Normalize
\begin{equation}
 U(X)=\frac{S(X)}{s_1X}
 =1+\sum_{j\ge1}u_jX^j,
 \qquad
 u_j=\frac{s_{j+1}}{s_1}.
 \label{eq:U-def}
\end{equation}
Then $q=S(X)^{1/2}=s_1^{1/2}X^{1/2}U(X)^{1/2}$.  Write
\begin{equation}
 T(z)=1+\sum_{r\ge1}t_rX^{r/2}.
 \label{eq:puiseux-def}
\end{equation}

\begin{resultbox}[title={All Puiseux coefficients}]
For every $r\ge1$,
\begin{equation}
 \boxed{
 t_r=-\!\sum_{\substack{1\le m\le r\\m\equiv r\ (\mathrm{mod}\,2)}}
 \omega_m s_1^{m/2}
 \Coeff X{(r-m)/2}{U(X)^{m/2}}.}
 \label{eq:tr-master}
\end{equation}
The remaining power coefficient is the finite Bell sum
\begin{equation}
 \Coeff XN{U(X)^{m/2}}
 =\frac1{N!}\sum_{k=0}^N\falling{m/2}k
 \pBell Nk(1!u_1,2!u_2,\ldots).
 \label{eq:tr-power-bell}
\end{equation}
Thus $t_r$ is a finite expression in $\omega_1,\ldots,\omega_r$ and
$\ell_1,\ldots,\ell_{(r+1)/2}$.
\end{resultbox}

\begin{proof}
Substitute $q=s_1^{1/2}X^{1/2}U^{1/2}$ into
$T=1-\sum_m\omega_mq^m$.  A term with index $m$ can contribute to
$X^{r/2}$ only if $r-m$ is a nonnegative even integer.  Extracting that
coefficient gives \cref{eq:tr-master}, and
\cref{eq:unit-power-bell} gives \cref{eq:tr-power-bell}.
\end{proof}

The first few coefficients, useful for hand checks, are
\begin{align}
 t_1={}&-\sqrt{s_1},\label{eq:t1}\\
 t_2={}&\frac{s_1}{3},\label{eq:t2}\\
 t_3={}&-\frac{11s_1^2+36s_2}{72\sqrt{s_1}},\label{eq:t3}\\
 t_4={}&\frac{43s_1^2+180s_2}{540},\label{eq:t4}\\
 t_5={}&-\frac{769s_1^4+3960s_1^2s_2+8640s_1s_3-2160s_2^2}
 {17280s_1^{3/2}},\label{eq:t5}\\
 t_6={}&\frac{442s_1^3+2709s_1s_2+5670s_3}{17010}.
 \label{eq:t6}
\end{align}

\section{Numerical Puiseux coefficients}
The first coefficients generated from \cref{eq:tr-master} are shown in
\cref{tab:puiseux}.  The displayed digits agree with the independent
calculations of Genitrini and continue directly to arbitrary order.

\begin{longtable}{r@{\qquad}r}
\caption{Coefficients in $T(z)=1+\sum_{r\ge1}t_r(1-z/\rho)^{r/2}$.}
\label{tab:puiseux}\\
\toprule
$r$ & $t_r$\\
\midrule
\endfirsthead
\toprule
$r$ & $t_r$\\
\midrule
\endhead
1  & $-1.559490020374640885542207396$\\
2  & $ \phantom{-}0.8106697078826992814381417462$\\
3  & $-0.2854870216128456053038513708$\\
4  & $ \phantom{-}0.1653723657120838934765135390$\\
5  & $-0.3424599704021542010422002117$\\
6  & $ \phantom{-}0.3174072259465285635680708874$\\
7  & $-0.1077788002916310091822310895$\\
8  & $ \phantom{-}0.06138495705583510468873149925$\\
9  & $-0.1952123835975564636127698720$\\
10 & $ \phantom{-}0.2059848312779074187361657348$\\
11 & $-0.05272470849819056231533821627$\\
12 & $ \phantom{-}0.01702656875495366944604539502$\\
13 & $-0.1523706243663253962860946421$\\
14 & $ \phantom{-}0.1737028832998504628939446698$\\
15 & $-0.01447370373952704485673111250$\\
16 & $-0.02189951761121556223267283747$\\
17 & $-0.1445471935709097054769461137$\\
18 & $ \phantom{-}0.1760771088850177790623322866$\\
\bottomrule
\end{longtable}

\begin{remark}[Analytic versus singular terms]
The even indices $t_{2j}X^j$ are analytic at $z=\rho$.  Any fixed analytic
monomial contributes no Taylor coefficients once $n>j$; the algebraic
coefficient asymptotics from the dominant point therefore depend only on the
odd coefficients $t_{2j+1}$.  The analytic part is not discarded globally: its
continuation participates in the exponentially smaller sectors studied in
Part~IV.
\end{remark}

\chapter{From Puiseux coefficients to the full asymptotic sequence}
\label{ch:transfer}

\section{Exact coefficient extraction for a singular monomial}
For $\beta\notin\N_0$,
\begin{equation}
 \Coeff zn{\left(1-\frac z\rho\right)^\beta}
 =\rho^{-n}\frac{\Gamma(n-\beta)}{\Gamma(-\beta)\Gamma(n+1)}.
 \label{eq:exact-binomial-transfer}
\end{equation}
The ratio of gamma functions has a complete inverse-power expansion.  Let
$B_m(x)$ denote the Bernoulli polynomials, with
$te^{xt}/(e^t-1)=\sum_{m\ge0}B_m(x)t^m/m!$.  Define
\begin{equation}
 d_m(\beta)=
 \frac{(-1)^{m+1}}{m(m+1)}
 \bigl(B_{m+1}(-\beta)-B_{m+1}(1)\bigr),
 \qquad m\ge1,
 \label{eq:d-beta}
\end{equation}
and let $g_r(\beta)$ be determined by
\begin{equation}
 \exp\!\left(\sum_{m\ge1}d_m(\beta)t^m\right)
 =\sum_{r\ge0}g_r(\beta)t^r.
 \label{eq:g-beta-def}
\end{equation}
Equivalently,
\begin{align}
 g_0(\beta)&=1,\label{eq:g0}\\
 g_r(\beta)&=\frac1r\sum_{m=1}^r m\,d_m(\beta)g_{r-m}(\beta),
 \label{eq:g-recurrence}\\
 g_r(\beta)&=\frac1{r!}\Bell r(1!d_1(\beta),2!d_2(\beta),\ldots,r!d_r(\beta)).
 \label{eq:g-bell}
\end{align}

\begin{theorem}[Gamma-ratio transfer coefficients]
\label{thm:gamma-ratio}
For fixed $\beta$ and $n\to\infty$ in a sector excluding the negative real axis,
\begin{equation}
 \frac{\Gamma(n-\beta)}{\Gamma(n+1)}
 \sim n^{-\beta-1}\sum_{r\ge0}\frac{g_r(\beta)}{n^r}.
 \label{eq:gamma-ratio-expansion}
\end{equation}
The coefficients are given by any of
\cref{eq:g-beta-def,eq:g-recurrence,eq:g-bell}.
\end{theorem}

\begin{proof}
The Bernoulli-polynomial form of Stirling's expansion gives
\[
 \log\frac{\Gamma(n-\beta)}{\Gamma(n+1)}
 =-(\beta+1)\log n+
 \sum_{m\ge1}\frac{d_m(\beta)}{n^m}.
\]
Exponentiating and using the complete Bell-polynomial identity proves the
claim.  The recurrence follows by differentiating
\cref{eq:g-beta-def} and comparing coefficients.
\end{proof}

For example,
\begin{equation}
 \bigl(g_r(1/2)\bigr)_{r\ge0}
 =1,\frac38,\frac{25}{128},\frac{105}{1024},
   \frac{1659}{32768},\frac{6237}{262144},\ldots .
 \label{eq:g-half}
\end{equation}

\section{All forward coefficients}
Combining the odd Puiseux terms with
\cref{eq:exact-binomial-transfer,eq:gamma-ratio-expansion} yields the desired
complete Poincar\'e sector.

\begin{theorem}[Full dominant expansion of A000081]
\label{thm:full-forward}
For every fixed $K\ge0$,
\begin{equation}
 a_n=\rho^{-n}n^{-3/2}
 \left(\sum_{k=0}^{K}\frac{A_k}{n^k}
 +O(n^{-K-1})\right),
 \qquad n\to\infty,
 \label{eq:forward-A}
\end{equation}
where
\begin{equation}
 \boxed{
 A_k=\sum_{j=0}^{k}
 \frac{t_{2j+1}}{\Gamma(-j-1/2)}
 g_{k-j}\!\left(j+\frac12\right).}
 \label{eq:Ak-master}
\end{equation}
Equivalently, writing $\Theta_k=\sqrt\pi A_k$,
\begin{equation}
 \boxed{
 \Theta_k=\sum_{j=0}^{k}
 t_{2j+1}g_{k-j}\!\left(j+\frac12\right)
 \frac{(2j+2)!}{(-4)^{j+1}(j+1)!}.}
 \label{eq:Theta-master}
\end{equation}
Every coefficient is therefore a finite expression in rational universal
constants, Bernoulli polynomials, Bell polynomials, and the disturbance jets
$\ell_1,\ldots,\ell_{k+1}$.
\end{theorem}

\begin{proof}
The term $t_{2j+1}X^{j+1/2}$ contributes
\[
 \rho^{-n}\frac{t_{2j+1}}{\Gamma(-j-1/2)}
 n^{-j-3/2}\sum_{r\ge0}\frac{g_r(j+1/2)}{n^r}.
\]
To obtain the coefficient of $n^{-k-3/2}$, set $r=k-j$ and sum over
$0\le j\le k$.  Finally use
\[
 \Gamma\!\left(-j-\frac12\right)
 =\frac{(-4)^{j+1}(j+1)!\sqrt\pi}{(2j+2)!}
\]
to obtain \cref{eq:Theta-master}.  Standard singularity analysis in a
$\Delta$-domain at the unique dominant singularity gives the stated remainder
at every fixed order \cite{FlajoletOdlyzko1990,Genitrini2016}.
\end{proof}

The first normalized transfer identities are particularly simple:
\begin{align}
 \Theta_0={}&-\frac12t_1,\label{eq:Theta0-t}\\
 \Theta_1={}&-\frac3{16}t_1+\frac34t_3,\label{eq:Theta1-t}\\
 \Theta_2={}&-\frac{25}{256}t_1+\frac{45}{32}t_3-\frac{15}{8}t_5,
 \label{eq:Theta2-t}\\
 \Theta_3={}&-\frac{105}{2048}t_1+\frac{1155}{512}t_3
              -\frac{525}{64}t_5+\frac{105}{16}t_7,
 \label{eq:Theta3-t}\\
 \Theta_4={}&-\frac{1659}{65536}t_1+\frac{14175}{4096}t_3
              -\frac{26775}{1024}t_5+\frac{6615}{128}t_7
              -\frac{945}{32}t_9.
 \label{eq:Theta4-t}
\end{align}

\section{The leading Otter constant}
Since $t_1=-\sqrt{2\ell_1}$,
\begin{equation}
 A_0=-\frac{t_1}{2\sqrt\pi}
 =\sqrt{\frac{\ell_1}{2\pi}}
 =\sqrt{\frac{1+\rho R'(\rho)}{2\pi}}.
 \label{eq:Otter-constant-formula}
\end{equation}
Numerically,
\begin{equation}
 A_0=0.43992401257102530404090339143454476479808540794012\ldots .
 \label{eq:A0-numeric}
\end{equation}
Thus the classical first-order form is
\begin{equation}
 a_n\sim
 0.4399240125710253040409\ldots\,
 (2.9557652856519949747148\ldots)^n n^{-3/2}.
 \label{eq:Otter-leading}
\end{equation}

\section{Numerical coefficient table}
It is convenient to display $\Theta_k=\sqrt\pi A_k$, so that
\begin{equation}
 a_n\sim\frac{\rho^{-n}}{\sqrt{\pi n^3}}
 \sum_{k\ge0}\frac{\Theta_k}{n^k}.
 \label{eq:forward-Theta}
\end{equation}

\begin{longtable}{r@{\qquad}r}
\caption{Forward coefficients in \cref{eq:forward-Theta}.}
\label{tab:forward-coeff}\\
\toprule
$k$ & $\Theta_k$\\
\midrule
\endfirsthead
\toprule
$k$ & $\Theta_k$\\
\midrule
\endhead
0  & $0.7797450101873204427711036979$\\
1  & $0.07828911261061096206127535867$\\
2  & $0.3929402676631860184743155979$\\
3  & $1.537879315978838091102498124$\\
4  & $8.200844090435596159219524856$\\
5  & $57.29291473494343787739929540$\\
6  & $503.0445050262735818247740217$\\
7  & $5359.600933884326033403142719$\\
8  & $67342.06920114653048340562550$\\
9  & $975425.4970695924732213025531$\\
10 & $15996932.49293173312960557140$\\
11 & $292822531.3353392231145687742$\\
12 & $5914523441.293932519117404924$\\
\bottomrule
\end{longtable}

Factoring out $A_0$ gives
\begin{equation}
 a_n\sim A_0\alpha^n n^{-3/2}F(n^{-1}),
 \qquad
 F(u)=1+\sum_{k\ge1}c_ku^k,
 \qquad c_k=\frac{A_k}{A_0}=\frac{\Theta_k}{\Theta_0}.
 \label{eq:F-forward}
\end{equation}
The first relative coefficients are
\begin{align}
 c_1={}&0.1004034800964014347637283944\ldots,\notag\\
 c_2={}&0.5039343151023035331021858170\ldots,\notag\\
 c_3={}&1.972284908382278469205675365\ldots,\notag\\
 c_4={}&10.51734090413157399273764936\ldots,\notag\\
 c_5={}&73.47647498402047106416577934\ldots,\notag\\
 c_6={}&645.1397552456612928153464923\ldots.
 \label{eq:c-numeric}
\end{align}

\section{Numerical validation and optimal truncation}
\Cref{tab:forward-errors} compares exact integers from
\cref{eq:exact-recurrence} with truncations after $n^{-K}$.  The entries are
signed relative errors $(a_n^{[K]}-a_n)/a_n$.

\begin{table}[htbp]
\centering
\caption{Signed relative errors of the forward expansion.}
\label{tab:forward-errors}
\begingroup
\scriptsize
\setlength{\tabcolsep}{3.1pt}
\begin{tabular}{r*{7}{r}}
\toprule
$n$ & $K=0$ & $K=1$ & $K=2$ & $K=4$ & $K=6$ & $K=8$ & $K=10$\\
\midrule
20  & $-6.58\,10^{-3}$ & $-1.59\,10^{-3}$ & $-3.42\,10^{-4}$ & $-3.18\,10^{-5}$ & $ 1.04\,10^{-6}$ & $ 9.73\,10^{-6}$ & $ 1.41\,10^{-5}$\\
50  & $-2.22\,10^{-3}$ & $-2.19\,10^{-4}$ & $-1.77\,10^{-5}$ & $-2.88\,10^{-7}$ & $-1.20\,10^{-8}$ & $-9.84\,10^{-10}$& $-1.36\,10^{-10}$\\
100 & $-1.06\,10^{-3}$ & $-5.24\,10^{-5}$ & $-2.08\,10^{-6}$ & $-8.06\,10^{-9}$ & $-7.88\,10^{-11}$& $-1.50\,10^{-12}$& $-4.73\,10^{-14}$\\
200 & $-5.15\,10^{-4}$ & $-1.28\,10^{-5}$ & $-2.53\,10^{-7}$ & $-2.40\,10^{-10}$& $-5.73\,10^{-13}$& $-2.66\,10^{-15}$& $-2.04\,10^{-17}$\\
500 & $-2.03\,10^{-4}$ & $-2.03\,10^{-6}$ & $-1.59\,10^{-8}$ & $-2.39\,10^{-12}$& $-9.02\,10^{-16}$& $-6.62\,10^{-19}$& $-8.01\,10^{-22}$\\
\bottomrule
\end{tabular}
\endgroup
\end{table}

At fixed $n$ the expansion is asymptotic, not convergent: the improvement stops
when the growing coefficients overcome the powers of $n$.  For example, at
$n=20$ the sixth-order truncation is better than the eighth- or tenth-order
truncation.  This late-term growth is the perturbative shadow of farther
singularities, which are made explicit in Part~IV.

\chapter{A reproducible coefficient pipeline}
\label{ch:algorithm}

\section{Dependency graph}
The coefficient calculation is triangular in the following sense.
\begin{center}
\begin{tikzpicture}[
 node distance=7mm and 10mm,
 box/.style={draw=deepolive,rounded corners,fill=softolive,align=center,
             minimum height=9mm,inner sep=2mm},
 arr/.style={-{Latex[length=2mm]},thick,draw=darkteal}]
\node[box] (a) {$a_1,\ldots,a_N$\\divisor recurrence};
\node[box,right=of a] (h) {$h_m$ and $R(z)$\\\cref{eq:h-coeff}};
\node[box,right=of h] (rho) {$\rho$ and jets $\ell_j$\\\cref{eq:rho-characterization,eq:ell-jets}};
\node[box,below=of rho] (s) {$s_j,U(X)$\\complete Bell};
\node[box,left=of s] (t) {$t_r$\\Lambert + partial Bell};
\node[box,left=of t] (A) {$A_k,c_k$\\gamma transfer};
\node[box,below=of A] (eta) {$\eta_k$\\logarithmic Bell transform};
\node[box,right=of eta] (inv) {inverse coefficients\\$D_k$ or $P_k(L)$};
\draw[arr] (a)--(h);
\draw[arr] (h)--(rho);
\draw[arr] (rho)--(s);
\draw[arr] (s)--(t);
\draw[arr] (t)--(A);
\draw[arr] (A)--(eta);
\draw[arr] (eta)--(inv);
\end{tikzpicture}
\end{center}
No high-order nonlinear system is solved.  At every stage, the coefficient of
order $r$ depends only on data of order at most $r$.

\section{Exact integer stage}
The following Python-like pseudocode realizes
\cref{eq:exact-recurrence}.  The array $b$ is updated as a divisor sieve.

\begin{lstlisting}[language=Python,caption={Exact generation of A000081.}]
def rooted_tree_numbers(N):
    a = [0] * (N + 1)
    b = [0] * (N + 1)   # b[m] = sum_{d | m} d*a[d]

    a[1] = 1
    for m in range(1, N + 1):
        b[m] += 1       # contribution of d = 1

    for n in range(2, N + 1):
        numerator = sum(b[j] * a[n-j] for j in range(1, n))
        assert numerator % (n - 1) == 0
        a[n] = numerator // (n - 1)

        for m in range(n, N + 1, n):
            b[m] += n * a[n]
    return a, b
\end{lstlisting}

The coefficients of $R$ are then exactly
\begin{equation}
 h_m=\frac{b_m-ma_m}{m}.
 \label{eq:algorithm-h}
\end{equation}
In fact $h_m$ is rational, though the defining proper-divisor sum often makes
its denominator cancellation transparent.

\section{Finding \texorpdfstring{$\rho$}{rho} and its jets}
For a target precision of $P$ decimal digits, choose $N$ so that
$(\sqrt\rho)^N$ is safely below $10^{-P}$, compute $a_1,\ldots,a_N$, and solve
\begin{equation}
 f(x)=\log x+1+\sum_{m=2}^{N}h_mx^m=0
 \label{eq:rho-truncated-equation}
\end{equation}
by Newton or safeguarded secant iteration.  The omitted tail is geometrically
small because $R$ is analytic in $|z|<\sqrt\rho$.  Once $\rho$ is known,
\cref{eq:R-derivative-sum} evaluates all required jets by rapidly convergent
sums.

\section{Formal-series stage}
For requested forward order $K$, it is enough to compute
\begin{enumerate}
\item $\ell_1,\ldots,\ell_{K+1}$;
\item $s_1,\ldots,s_{K+1}$ by \cref{eq:sj-bell};
\item $\omega_1,\ldots,\omega_{2K+1}$ by \cref{eq:omega-master};
\item $t_1,t_3,\ldots,t_{2K+1}$ by \cref{eq:tr-master};
\item $g_r(j+1/2)$ for $j+r\le K$ by \cref{eq:g-recurrence};
\item $A_0,\ldots,A_K$ by \cref{eq:Ak-master}.
\end{enumerate}
Naive polynomial arithmetic costs $O(K^3)$ scalar operations.  Bell recurrences,
FFT multiplication, and relaxed power-series methods reduce this when hundreds
or thousands of coefficients are required.  For the orders normally useful in
numerical asymptotics, the naive triangular implementation is simpler and more
reliable.

\begin{warningbox}[title={Precision discipline}]
The asymptotic coefficients eventually grow rapidly.  A calculation of
$A_K$ with $d$ reliable digits may require substantially more than $d$ working
digits because \cref{eq:Ak-master} combines alternating contributions from
several Puiseux coefficients.  The calculation underlying the tables used 90--120
decimal working digits and exact rational arithmetic for the universal
coefficients and gamma-ratio coefficients.
\end{warningbox}

\part{The Asymptotic Inverse}

\chapter{What it means to invert a sequence}
\label{ch:inverse-meaning}

\section{Three inverse objects}
A sequence is not literally a real-variable function, and A000081 begins with
$a_1=a_2=1$.  The phrase ``inverse of the sequence'' can therefore denote three
different objects.

\begin{definition}[Smooth asymptotic inverse]
Let
\begin{equation}
 \mathcal A(x)=A_0\alpha^x x^{-3/2}F(x^{-1})
 \label{eq:smooth-interpolant}
\end{equation}
be the formal asymptotic interpolation obtained from
\cref{eq:F-forward}.  Its asymptotic inverse $\nu(y)$ is the formal solution of
$\mathcal A(\nu(y))=y$ with $\nu(y)\to+\infty$.
\end{definition}

\begin{definition}[Generalized integer inverse]
For $y>1$, define
\begin{equation}
 N_*(y)=\min\{n\ge2:a_n\ge y\}.
 \label{eq:integer-inverse}
\end{equation}
This is a right-continuous staircase.
\end{definition}

\begin{definition}[Exact interpolated inverse]
Choose a monotone interpolation $\mathcal A_{\rm ex}(x)$ satisfying
$\mathcal A_{\rm ex}(n)=a_n$; for example, interpolate $\log a_n$ linearly on
each interval $[n,n+1]$.  Let $\nu_{\rm ex}$ be its ordinary inverse.
\end{definition}

The smooth inverse is canonical at the level of asymptotic series.  The exact
interpolated inverse depends on the chosen interpolation between integers, and
the generalized integer inverse retains an order-one rounding oscillation.
These distinctions do not affect the coefficients derived below, but they do
affect claims of pointwise accuracy smaller than one.

\section{The unavoidable rounding layer}
For nonintegral $x$,
\begin{equation}
 \lceil x\rceil
 =x+\frac12+\frac1\pi\sum_{m\ge1}\frac{\sin(2\pi mx)}m,
 \label{eq:ceiling-fourier}
\end{equation}
with the usual midpoint interpretation at integers.  Thus, once an exact
monotone interpolation has been chosen,
\begin{equation}
 N_*(y)=\lceil\nu_{\rm ex}(y)\rceil.
 \label{eq:integer-via-exact}
\end{equation}
The Fourier series in \cref{eq:ceiling-fourier} is not a small correction: it
is an $O(1)$ periodic layer.  Consequently no smooth function can approximate
$N_*(y)$ uniformly with error $o(1)$ for all real $y$.  At the special values
$y=a_n$, however, the smooth inverse satisfies
$\nu(a_n)=n+o(1)$, and successive asymptotic orders make this error arbitrarily
small before the divergent late-term regime is reached.

\begin{warningbox}
An expansion such as $N_*(y)=\nu(y)+o(1)$ is false uniformly in $y$.  The correct
statement is either about the smooth inverse $\nu$, about values $y=a_n$, or
about a rounded exact interpolation.  The integer staircase and the analytic
inverse must not be silently identified.
\end{warningbox}

\chapter{Normalization of the inverse problem}
\label{ch:inverse-normalization}

\section{Logarithmic phase}
Set
\begin{equation}
 p=\frac32,
 \qquad
 C=A_0,
 \qquad
 \lambda=\log\alpha=-\log\rho.
 \label{eq:p-C-lambda}
\end{equation}
Starting from
$y=C\e^{\lambda x}x^{-p}F(x^{-1})$, introduce
\begin{equation}
 X=\lambda x,
 \qquad
 Z=\log\frac{y}{C\lambda^p}.
 \label{eq:XZ}
\end{equation}
Then
\begin{equation}
 Z=X-p\log X+K(X^{-1}),
 \label{eq:normalized-inverse-equation}
\end{equation}
where
\begin{equation}
 K(v)=\log F(\lambda v)=\sum_{k\ge1}\eta_kv^k.
 \label{eq:K-eta}
\end{equation}
The logarithm is the key simplification: the multiplicative forward correction
becomes an additive inverse-power phase.

\section{General coefficient formula for the phase}
By \cref{eq:log-bell},
\begin{equation}
 \boxed{
 \eta_n=\frac1{n!}\sum_{k=1}^n(-1)^{k-1}(k-1)!
 \pBell nk(1!c_1\lambda,2!c_2\lambda^2,\ldots).}
 \label{eq:eta-master}
\end{equation}
In particular,
\begin{align}
 \eta_1&=\lambda c_1,\notag\\
 \eta_2&=\lambda^2\left(c_2-\frac12c_1^2\right),\notag\\
 \eta_3&=\lambda^3\left(c_3-c_1c_2+\frac13c_1^3\right),\notag\\
 \eta_4&=\lambda^4\left(c_4-c_1c_3-\frac12c_2^2+c_1^2c_2-\frac14c_1^4\right).
 \label{eq:eta-first-symbolic}
\end{align}
Numerically,
\begin{align}
 \eta_1={}&0.1088130343442745145167554784\ldots,\notag\\
 \eta_2={}&0.5859660997717814776889432009\ldots,\notag\\
 \eta_3={}&2.446558573456779649881479860\ldots,\notag\\
 \eta_4={}&14.06753235112498242303023266\ldots,\notag\\
 \eta_5={}&106.8546611534523691271119951\ldots,\notag\\
 \eta_6={}&1022.176467223247872767588518\ldots.
 \label{eq:eta-numeric}
\end{align}
Also
\begin{equation}
 C\lambda^p
 =0.4963361416595997474727369464252249134814\ldots .
 \label{eq:Clambdap}
\end{equation}

\chapter{\texorpdfstring{Lambert--$W$}{Lambert--W} core and an ordinary correction series}
\label{ch:inverse-lambert}

\section{Exact summation of the logarithmic skeleton}
First omit $K$ and solve
\begin{equation}
 X_0-p\log X_0=Z.
 \label{eq:Lambert-core-equation}
\end{equation}
The large positive solution is
\begin{equation}
 \boxed{
 X_0=-pW_{-1}\!\left(-\frac1p\e^{-Z/p}\right).}
 \label{eq:X0-W}
\end{equation}
Indeed, with $r=X_0/p$,
$r\e^{-r}=\e^{-Z/p}/p$, and the large solution corresponds to the branch
$W_{-1}$.  This single expression resums the complete tower
$Z+p\log Z+p^2\log Z/Z+\cdots$ generated by the power-law prefactor
$x^{-p}$.

\section{Correction about the exact core}
Write
\begin{equation}
 X=X_0+\delta.
 \label{eq:X-delta}
\end{equation}
Subtracting \cref{eq:Lambert-core-equation} from
\cref{eq:normalized-inverse-equation} gives
\begin{equation}
 \delta=p\log\left(1+\frac{\delta}{X_0}\right)
 -K\!\left(\frac1{X_0+\delta}\right).
 \label{eq:delta-fixed-point}
\end{equation}
Set $\varepsilon=X_0^{-1}$ and define
\begin{equation}
 \Phi_\varepsilon(u)
 =p\log(1+\varepsilon u)
  -K\!\left(\frac{\varepsilon}{1+\varepsilon u}\right).
 \label{eq:Phi-epsilon}
\end{equation}
Then $\delta=\Phi_\varepsilon(\delta)$ and
$\Phi_\varepsilon=O(\varepsilon)$.

\begin{theorem}[All inverse corrections about the Lambert core]
\label{thm:delta-Lagrange}
As a formal series in $X_0^{-1}$,
\begin{equation}
 \boxed{
 \delta=\sum_{m\ge1}\frac1{m!}
 \left.\frac{\dd^{m-1}}{\dd u^{m-1}}
 \Phi_\varepsilon(u)^m\right|_{u=0},
 \qquad \varepsilon=X_0^{-1}.}
 \label{eq:delta-Lagrange}
\end{equation}
Equivalently,
\begin{equation}
 \delta\sim\sum_{r\ge1}\frac{D_r}{X_0^r},
 \label{eq:delta-D}
\end{equation}
where the completely explicit coefficient extractor is
\begin{equation}
 \boxed{
 D_r=[\varepsilon^r]
 \sum_{m=1}^{r}\frac1{m!}
 \left.\frac{\dd^{m-1}}{\dd u^{m-1}}
 \left[p\log(1+\varepsilon u)
 -K\!\left(\frac{\varepsilon}{1+\varepsilon u}\right)\right]^m
 \right|_{u=0}.}
 \label{eq:Dr-master}
\end{equation}
Finally,
\begin{equation}
 \boxed{
 \nu(y)\sim\frac1\lambda\left(
 X_0+\sum_{r\ge1}\frac{D_r}{X_0^r}\right),
 \quad
 X_0=-pW_{-1}\!\left(-\frac1p
 \left(\frac{C\lambda^p}{y}\right)^{1/p}\right).}
 \label{eq:inverse-W-master}
\end{equation}
\end{theorem}

\begin{proof}
Introduce an auxiliary marker $t$ and solve
$\delta=t\Phi_\varepsilon(\delta)$.  Lagrange--B\"urmann gives
\[
 \delta(t)=\sum_{m\ge1}\frac{t^m}{m!}
 \left.\frac{\dd^{m-1}}{\dd u^{m-1}}
 \Phi_\varepsilon(u)^m\right|_{u=0}.
\]
Since $\Phi_\varepsilon=O(\varepsilon)$, the coefficient of
$\varepsilon^r$ receives contributions only from $m\le r$, so setting $t=1$
is formally legitimate and gives \cref{eq:Dr-master}.  Division by $\lambda$
returns from $X$ to $x$.
\end{proof}

\section{First correction coefficients}
The first $D_r$ are
\begin{align}
 D_1={}&-\eta_1,\label{eq:D1}\\
 D_2={}&-p\eta_1-\eta_2,\label{eq:D2}\\
 D_3={}&-p^2\eta_1-p\eta_2-\eta_1^2-\eta_3,\label{eq:D3}\\
 D_4={}&-p^3\eta_1-p^2\eta_2-\frac52p\eta_1^2
         -p\eta_3-3\eta_1\eta_2-\eta_4,\label{eq:D4}\\
 D_5={}&-p^4\eta_1-p^3\eta_2-\frac92p^2\eta_1^2-p^2\eta_3
         -7p\eta_1\eta_2-p\eta_4\notag\\
      &\hspace{12mm}-2\eta_1^3-4\eta_1\eta_3-2\eta_2^2-\eta_5.
 \label{eq:D5}
\end{align}
For A000081 these are
\begin{align}
 D_1={}&-0.1088130343442745145167554784\ldots,\notag\\
 D_2={}&-0.7491856512881932494640764185\ldots,\notag\\
 D_3={}&-3.582177326832277789103009844\ldots,\notag\\
 D_4={}&-19.65872121138776003707556394\ldots,\notag\\
 D_5={}&-138.5327478518560861313634930\ldots,\notag\\
 D_6={}&-1248.979326371692854496506550\ldots,\notag\\
 D_7={}&-13905.40884474814835419890912\ldots,\notag\\
 D_8={}&-184719.1911265080904894459391\ldots.
 \label{eq:D-numeric}
\end{align}
The coefficients grow, as expected for an asymptotic series; $W_{-1}$ has
already removed the dominant logarithmic combinatorics, but it cannot remove
the farther singularities encoded in the forward coefficients.

\chapter{Polynomial--logarithmic inverse transseries}
\label{ch:polylog-inverse}

\section{Direct expansion without a special function}
Some applications prefer an expansion involving only elementary logarithms.
Let
\begin{equation}
 L=\log Z,
 \qquad
 X=Z+U,
 \qquad
 U=\sum_{n\ge0}P_n(L)Z^{-n}.
 \label{eq:P-ansatz}
\end{equation}
With $\varepsilon=Z^{-1}$, substitution into
\cref{eq:normalized-inverse-equation} gives a single generating equation for
all coefficient polynomials:
\begin{equation}
 \boxed{
 P(\varepsilon,L)=pL+p\log\bigl(1+\varepsilon P(\varepsilon,L)\bigr)
 -K\!\left(\frac{\varepsilon}
 {1+\varepsilon P(\varepsilon,L)}\right),}
 \label{eq:P-functional}
\end{equation}
where
$P(\varepsilon,L)=\sum_{n\ge0}P_n(L)\varepsilon^n$.

\section{A closed triangular recurrence}
Define convolution polynomials
\begin{equation}
 C_{m,r}(L)=[\varepsilon^m]P(\varepsilon,L)^r
 =\sum_{j_1+\cdots+j_r=m}P_{j_1}(L)\cdots P_{j_r}(L),
 \label{eq:Cmr}
\end{equation}
with $C_{0,0}=1$ and $C_{m,0}=0$ for $m>0$.  Expanding the logarithm and the
negative powers in \cref{eq:P-functional} yields the following all-order
recurrence.

\begin{theorem}[Coefficient polynomials of the direct inverse]
\label{thm:P-recurrence}
One has $P_0(L)=pL$, and for $n\ge1$,
\begin{equation}
 \boxed{\begin{aligned}
 P_n(L)={}&p\sum_{r=1}^{n}\frac{(-1)^{r+1}}r C_{n-r,r}(L)\\
 &-\sum_{k=1}^{n}\eta_k
   \sum_{r=0}^{n-k}(-1)^r\binom{k+r-1}{r}
   C_{n-k-r,r}(L)
 \end{aligned}}
 \label{eq:Pn-master}
\end{equation}
Moreover $P_n$ is a polynomial of degree $n$ for $n\ge1$, and
\begin{equation}
 \boxed{
 \nu(y)\sim\frac1\lambda\left[
 Z+p\log Z+\sum_{n\ge1}\frac{P_n(\log Z)}{Z^n}\right],
 \quad Z=\log\frac{y}{C\lambda^p}.}
 \label{eq:inverse-polylog-master}
\end{equation}
\end{theorem}

\begin{proof}
Use
\[
 \log(1+\varepsilon P)
 =\sum_{r\ge1}\frac{(-1)^{r+1}}r\varepsilon^rP^r
\]
and
\[
 K\!\left(\frac{\varepsilon}{1+\varepsilon P}\right)
 =\sum_{k\ge1}\eta_k\varepsilon^k
  \sum_{r\ge0}(-1)^r\binom{k+r-1}{r}\varepsilon^rP^r.
\]
Extracting $\varepsilon^n$ gives \cref{eq:Pn-master}.  The right side at
order $n$ uses only $P_0,\ldots,P_{n-1}$, so the recurrence is triangular.
Induction on $n$ proves the degree bound, and substitution establishes the
asymptotic expansion.
\end{proof}

\section{First coefficient polynomials}
The first five are
\begin{align}
 P_0(L)={}&pL,\label{eq:P0}\\
 P_1(L)={}&p^2L-\eta_1,\label{eq:P1}\\
 P_2(L)={}&-\frac12p^3L^2+(p^3+p\eta_1)L-p\eta_1-\eta_2,
 \label{eq:P2}\\
 P_3(L)={}&\frac13p^4L^3-
 \left(\frac32p^4+p^2\eta_1\right)L^2\notag\\
 &+\left(p^4+3p^2\eta_1+2p\eta_2\right)L
 -p^2\eta_1-p\eta_2-\eta_1^2-\eta_3,
 \label{eq:P3}\\
 P_4(L)={}&-\frac14p^5L^4
 +\left(\frac{11}{6}p^5+p^3\eta_1\right)L^3\notag\\
 &-\left(3p^5+\frac{11}{2}p^3\eta_1+3p^2\eta_2\right)L^2\notag\\
 &+\left(p^5+6p^3\eta_1+5p^2\eta_2
          +3p\eta_1^2+3p\eta_3\right)L\notag\\
 &-p^3\eta_1-p^2\eta_2-\frac52p\eta_1^2-p\eta_3
   -3\eta_1\eta_2-\eta_4.
 \label{eq:P4}
\end{align}
Notice that
\begin{equation}
 D_n=P_n(0).
 \label{eq:D-P0}
\end{equation}
This is not accidental: the Lambert core sums exactly the part of the direct
recurrence generated by $pL$; the remaining constants obey the same fixed-point
calculus with $L=0$.

For $p=3/2$ and the A000081 phase coefficients,
\begin{align}
 P_0(L)={}&1.5L,\notag\\
 P_1(L)={}&2.25L-0.10881303434427451452\ldots,\notag\\
 P_2(L)={}&-1.6875L^2+3.5382195515164117718\ldots L\notag\\
          &\quad-0.74918565128819324946\ldots,\notag\\
 P_3(L)={}&1.6875L^3-7.8385793272746176577\ldots L^2\notag\\
          &\quad+7.5548862811391974061\ldots L
            -3.5821773268322777891\ldots,\notag\\
 P_4(L)={}&-1.8984375L^4+14.289118990911926486\ldots L^3\notag\\
          &\quad-28.756363123475120650\ldots L^2\notag\\
          &\quad+27.452127392454046160\ldots L
            -19.658721211387760037\ldots .
 \label{eq:P-numeric}
\end{align}

\section{Elementary leading form}
Expanding only to constant accuracy in terms of $Y=\log y$ gives
\begin{equation}
 \nu(y)=\frac{\log y}{\lambda}
 +\frac{p}{\lambda}\log\log y
 -\frac1\lambda\log(C\lambda^p)+o(1).
 \label{eq:inverse-leading-elementary}
\end{equation}
Numerically,
\begin{align}
 \nu(y)={}&0.92271556161860152480\ldots\log y\notag\\
 &+1.38407334242790228720\ldots\log\log y\notag\\
 &+0.64636398270020880729\ldots+o(1).
 \label{eq:inverse-leading-numeric}
\end{align}
The constant in this coarser form depends on the exact normalization of $Z$;
using \cref{eq:inverse-polylog-master} is preferable whenever more than the
first two scales are needed.

\chapter{Accuracy, implementation, and integer recovery}
\label{ch:inverse-practical}

\section{Comparison at exact sequence values}
For $y=a_n$, define the order-$K$ Lambert-core approximation
\begin{equation}
 \nu^{[K]}(a_n)=\frac1\lambda
 \left(X_0+\sum_{r=1}^{K}\frac{D_r}{X_0^r}\right).
 \label{eq:inverse-truncation}
\end{equation}
\Cref{tab:inverse-errors} gives the signed index error
$\nu^{[K]}(a_n)-n$.

\begin{table}[htbp]
\centering
\caption{Signed error of the asymptotic inverse at $y=a_n$.}
\label{tab:inverse-errors}
\begingroup
\scriptsize
\setlength{\tabcolsep}{3.1pt}
\begin{tabular}{r*{7}{r}}
\toprule
$n$ & $K=0$ & $K=1$ & $K=2$ & $K=4$ & $K=6$ & $K=8$ & $K=10$\\
\midrule
20  & $6.55\,10^{-3}$ & $1.91\,10^{-3}$ & $4.44\,10^{-4}$ & $3.76\,10^{-5}$ & $-2.02\,10^{-7}$ & $-9.39\,10^{-6}$ & $-1.39\,10^{-5}$\\
50  & $2.11\,10^{-3}$ & $2.59\,10^{-4}$ & $2.32\,10^{-5}$ & $3.32\,10^{-7}$ & $1.26\,10^{-8}$ & $9.98\,10^{-10}$& $1.35\,10^{-10}$\\
100 & $9.88\,10^{-4}$ & $6.16\,10^{-5}$ & $2.74\,10^{-6}$ & $9.34\,10^{-9}$ & $8.35\,10^{-11}$& $1.53\,10^{-12}$& $4.71\,10^{-14}$\\
200 & $4.78\,10^{-4}$ & $1.50\,10^{-5}$ & $3.33\,10^{-7}$ & $2.79\,10^{-10}$& $6.09\,10^{-13}$& $2.71\,10^{-15}$& $2.03\,10^{-17}$\\
500 & $1.88\,10^{-4}$ & $2.38\,10^{-6}$ & $2.10\,10^{-8}$ & $2.78\,10^{-12}$& $9.59\,10^{-16}$& $6.75\,10^{-19}$& $7.98\,10^{-22}$\\
\bottomrule
\end{tabular}
\endgroup
\end{table}

The inverse error is essentially the forward relative error divided by the
logarithmic slope $\lambda-p/n+O(n^{-2})$.  This explains the close numerical
parallel between \cref{tab:forward-errors,tab:inverse-errors}.

\section{A robust numerical procedure}
For a large positive target $y$, a stable implementation is:
\begin{algorithm}[Asymptotic inversion of A000081]
\leavevmode
\begin{enumerate}
\item Compute $Z=\log(y/(C\lambda^{3/2}))$ using a logarithmic representation of
      $y$ if $y$ is too large for ordinary floating point.
\item Evaluate
\[
 X_0=-\frac32W_{-1}\!\left(-\frac23\e^{-2Z/3}\right).
\]
\item Choose a truncation order $K$ and form
$X=X_0+\sum_{r=1}^{K}D_rX_0^{-r}$.
\item Return $x=X/\lambda$.  If an integer index is desired, test the two or
      three integers nearest to $x$ by exact recurrence values, arbitrary
      precision logarithms, or a certified forward bound.
\end{enumerate}
\end{algorithm}

A single Newton correction can be applied to a chosen smooth truncation:
\begin{equation}
 x_{\rm new}=x-
 \frac{\log\mathcal A_K(x)-\log y}
 {\lambda-p/x+\frac{\dd}{\dd x}\log F_K(x^{-1})}.
 \label{eq:inverse-newton}
\end{equation}
Working with logarithms avoids overflow and loss of relative precision.

\section{Rounding certification}
Because $a_{n+1}/a_n\to\alpha>1$, adjacent sequence values eventually differ by
a fixed multiplicative factor.  Hence an index estimate with absolute error
less than $1/2$ identifies the nearest integer at targets $y=a_n$.  For a
general target $y$, certification of $N_*(y)$ requires checking
\begin{equation}
 a_{m-1}<y\le a_m.
 \label{eq:inverse-certification}
\end{equation}
The asymptotic inverse reduces this check to a constant-size neighbourhood of
$m\approx\nu(y)$; it does not replace the endpoint convention in the definition
of the staircase.

\part{Beyond All Orders: The Singularity Transseries}

\chapter{The recursive singularity atlas}
\label{ch:atlas}

\section{Why one exponential scale is not the whole story}
The expansion in \cref{eq:forward-A} is complete in inverse powers of $n$ at
the dominant exponential scale $\rho^{-n}$.  It does not contain terms such as
$\sigma^{-n}$ with $|\sigma|>\rho$, because those are smaller than every power
of $1/n$ relative to $\rho^{-n}$.  Such terms arise from farther singularities
of analytic continuation.  For the P\'olya-tree equation there are infinitely
many of them, and they accumulate at the unit circle, which is a natural
boundary \cite{DrmotaGittenberger2010}.

The functional equation exposes two mechanisms.
\begin{enumerate}
\item \textbf{Critical Lambert singularities.}  On a branch where $R$ is
      analytic, a point $\sigma$ satisfying
      $\zeta(\sigma)=1/e$ is a square-root singularity whenever
      $\zeta'(\sigma)\ne0$.
\item \textbf{Inherited singularities.}  If $T$ is singular at $\tau$ and
      $\sigma^k=\tau$ for some $k\ge2$, then the term $T(z^k)/k$ can make
      $R(z)$, hence $T(z)$, singular at $\sigma$.  Repeated root lifts produce
      singularities whose moduli approach one and whose arguments become
      dense.
\end{enumerate}
Different paths can reach different determinations of the inner functions and
of the outer Lambert function.  The appropriate global object is therefore the
Riemann surface of the germ $T$ at zero, not a single-valued function on the
whole punctured unit disk.

\section{Local closure of the Puiseux calculus}
Fix a continuation sheet and a point $\sigma\ne0$.  Choose a local uniformizer
$v$ such that
\begin{equation}
 z=\sigma(1-v^M)
 \label{eq:local-uniformizer}
\end{equation}
for some positive integer $M$, and suppose all inner germs needed in $R(z)$
have Puiseux expansions in $v$.  Algebraic addition, composition with $z^k$,
exponentiation, and multiplication preserve the ring $\C[[v]]$.  Thus
\begin{equation}
 \zeta(v)=\zeta_0+\sum_{m\ge1}\zeta_mv^m.
 \label{eq:zeta-local-v}
\end{equation}
There are two cases.

\subsection{Regular outer composition}
Let $C_b$ denote the branch of the Cayley function reached on the chosen sheet,
and assume
\begin{equation}
 y_0=C_b(\zeta_0)\ne1.
 \label{eq:regular-outer-condition}
\end{equation}
Then $C_b$ is analytic at $\zeta_0$, and
\begin{equation}
 T(v)=y_0+\sum_{m\ge1}y_mv^m,
 \qquad
 \boxed{
 y_m=\frac1{m!}\sum_{k=1}^m C_b^{(k)}(\zeta_0)
 \pBell mk(1!\zeta_1,2!\zeta_2,\ldots).}
 \label{eq:regular-local-composition}
\end{equation}
The derivatives are explicit functions of $y_0$.  Since
$u=y\e^{-y}$,
\begin{equation}
 C_b'(u)=\frac{\e^y}{1-y},
 \qquad
 C_b^{(k)}(u)=
 \left[\left(\frac{\e^y}{1-y}\frac{\dd}{\dd y}\right)^{k-1}
 \frac{\e^y}{1-y}\right]_{y=C_b(u)}.
 \label{eq:C-derivative-operator}
\end{equation}

\subsection{Critical outer composition}
If $\zeta_0=1/e$ and the reached branch has $y_0=1$, define
\begin{equation}
 Q(v)=\sqrt{2(1-e\zeta(v))}.
 \label{eq:Q-local}
\end{equation}
Then the universal formula of \cref{thm:omega} gives
\begin{equation}
 \boxed{
 T(v)=1-\sum_{m\ge1}\omega_m Q(v)^m.}
 \label{eq:critical-local-composition}
\end{equation}
If $1-e\zeta(v)$ vanishes first at order $r$, then $Q$ begins with
$v^{r/2}$; the ramification index doubles when $r$ is odd.  This is the only
new ramification mechanism.  Consequently the complete singularity atlas can
be generated recursively by ordinary formal power-series arithmetic and the
same universal sequence $(\omega_m)$.

\begin{theorem}[Recursive completeness of the local calculus]
\label{thm:local-completeness}
Starting from the dominant germ at $\rho$, every isolated singular germ reached
by finitely many applications of the root-lift and critical mechanisms admits a
computable Puiseux expansion.  At each fixed local order, its coefficients are
finite expressions built from
\begin{enumerate}
\item previously computed local coefficients of inner germs;
\item complete and partial Bell polynomials;
\item the rational Lambert coefficients $\omega_m$;
\item derivatives of a regular Cayley branch as in
      \cref{eq:C-derivative-operator}.
\end{enumerate}
\end{theorem}

\begin{proof}
Proceed by induction on the depth of the singularity in the root-lift graph.
At depth zero, \cref{eq:tr-master} gives the dominant germ.  At the induction
step, substitution $z\mapsto z^k$ converts every already known Puiseux germ to
a Puiseux series in a common local uniformizer.  The sum defining $R$, its
exponential, and multiplication by $z$ preserve that ring.  If the outer Cayley
branch is regular, apply \cref{eq:regular-local-composition}; if it is critical,
apply \cref{eq:critical-local-composition}.  Both formulae determine every
coefficient at finite order.
\end{proof}

\section{Contour-resolved coefficient transseries}
Let a chosen deformation of the Cauchy coefficient contour encounter isolated
singularities $\sigma$ on specified sheets.  Suppose the local singular part is
\begin{equation}
 T_\sigma(z)\asympseries
 \sum_{q\in E_\sigma}c_{\sigma,q}
 \left(1-\frac z\sigma\right)^q,
 \label{eq:local-singular-general}
\end{equation}
where $E_\sigma\subset\mathbb Q_{\ge0}$ and analytic exponents
$q\in\N_0$ may be omitted from the local Hankel contribution.  Let
$\Stokes_\sigma$ be the contour intersection or Stokes multiplier associated
with that local branch.

\begin{resultbox}[title={Full contour-resolved forward transseries}]
Under the usual finite-radius continuation and remainder hypotheses for the
chosen contour,
\begin{align}
 a_n\asympseries{}&
 \sum_{\sigma}\Stokes_\sigma\,\sigma^{-n}
 \sum_{q\in E_\sigma\setminus\N_0}
 \frac{c_{\sigma,q}}{\Gamma(-q)}n^{-q-1}
 \sum_{r\ge0}\frac{g_r(q)}{n^r}.
 \label{eq:global-forward-transseries}
\end{align}
Relative to the dominant sector, define the action
\begin{equation}
 \Action_\sigma=\Log\frac{\sigma}{\rho}.
 \label{eq:action}
\end{equation}
Then $\sigma^{-n}=\rho^{-n}\e^{-n\Action_\sigma}$, so
\cref{eq:global-forward-transseries} is a genuine exponential transseries.
Conjugate singularities combine into real oscillatory sectors.
\end{resultbox}

\begin{proof}[Coefficient-contour argument]
Deform the Cauchy contour from a small circle to a finite union of local Hankel
contours about the encountered singularities plus an outer remainder contour.
Apply \cref{eq:exact-binomial-transfer} termwise to each truncated local
Puiseux expansion and then use \cref{thm:gamma-ratio}.  The outer contour gives
the radius-dependent exponential remainder.  Passing through a branch cut or
to a different sheet changes the homology coefficients of the local Hankel
cycles; these are precisely the multipliers $\Stokes_\sigma$.
\end{proof}

\begin{remark}[Projective meaning of the infinite sum]
Because singularities accumulate at $|z|=1$, the sum over $\sigma$ is not meant
as an absolutely convergent unordered numerical series.  For every fixed
continuation radius $R<1$, only the singularities met before the $R$-contour are
included, and the rest is absorbed into the outer remainder.  The global
transseries is the compatible family of these finite-radius expansions as the
contour is enlarged.  This is the natural notion of fullness in the presence of
a natural boundary.
\end{remark}

\chapter{The first inherited sector at \texorpdfstring{$-\sqrt\rho$}{-sqrt(rho)}}
\label{ch:negative-sector}

\section{A singularity accessible along the negative axis}
Let
\begin{equation}
 \sigma_-=-\sqrt\rho
 =-0.5816544136333942538100577816448348312833\ldots .
 \label{eq:sigma-minus}
\end{equation}
The germ at the origin can be continued along the negative real axis without
meeting the dominant positive singularity.  At $z=\sigma_-$, the term
$T(z^2)/2$ in $R(z)$ reaches $T(\rho)/2$ and inherits the dominant square-root
branch.  All terms $T(z^k)$ with $k\ge3$ are evaluated strictly inside the
original disk of convergence because $|\sigma_-|^k<\rho$.

Choose
\begin{equation}
 z=\sigma_-(1-v^2),
 \qquad v=\left(1-\frac z{\sigma_-}\right)^{1/2}.
 \label{eq:v-minus}
\end{equation}
Then
\[
 z^2=\rho(1-v^2)^2,
 \qquad
 1-\frac{z^2}{\rho}=v^2(2-v^2).
\]
Using the dominant coefficients $t_r$, the complete local disturbance is
\begin{align}
 R_-(v)={}&\frac12\left[
 1+\sum_{r\ge1}t_rv^r(2-v^2)^{r/2}\right]\notag\\
 &+\sum_{k\ge3}\frac1k\sum_{n\ge1}
 a_n\sigma_-^{kn}(1-v^2)^{kn}.
 \label{eq:Rminus-exact}
\end{align}
Both lines are convergent as local Puiseux/power series near $v=0$.

\section{General coefficients of the inherited disturbance}
Write $R_-(v)=\sum_{m\ge0}r_m^-v^m$ and put $t_0=1$.  For
$j=(m-r)/2$, direct coefficient extraction from
\cref{eq:Rminus-exact} gives
\begin{align}
 r_m^-={}&\frac12
 \sum_{\substack{0\le r\le m\\r\equiv m\ (\mathrm{mod}\,2)}}
 t_r(-1)^j2^{r/2-j}\binom{r/2}{j}\notag\\
 &+\mathbf 1_{2\mid m}(-1)^{m/2}
 \sum_{k\ge3}\frac1k\sum_{n\ge1}a_n\sigma_-^{kn}
 \binom{kn}{m/2}.
 \label{eq:rminus-master}
\end{align}
This is an explicit all-order formula.  The double sum in the second line
converges geometrically, already at $k=3$.

Define
\begin{equation}
 \zeta_-(v)=\sigma_-(1-v^2)\e^{R_-(v)}
 =\sum_{m\ge0}z_m^-v^m.
 \label{eq:zetaminus-series}
\end{equation}
If
\begin{equation}
 E_m^-=\frac{\e^{r_0^-}}{m!}
 \Bell m(1!r_1^-,2!r_2^-,\ldots,m!r_m^-),
 \qquad E_m^-=0\quad(m<0),
 \label{eq:Eminus}
\end{equation}
then
\begin{equation}
 z_m^-=\sigma_-(E_m^--E_{m-2}^-).
 \label{eq:zminus-bell}
\end{equation}
Finally let
\begin{equation}
 y_-=C_0(z_0^-)=-W_0(-z_0^-),
 \qquad
 T(\sigma_-(1-v^2))=y_-+\sum_{m\ge1}b_m^-v^m.
 \label{eq:bminus-def}
\end{equation}
The outer Cayley branch is regular because $y_-\ne1$, so
\begin{equation}
 \boxed{
 b_m^-=\frac1{m!}\sum_{k=1}^m C_0^{(k)}(z_0^-)
 \pBell mk(1!z_1^-,2!z_2^-,\ldots).}
 \label{eq:bminus-master}
\end{equation}
Together,
\cref{eq:rminus-master,eq:zminus-bell,eq:bminus-master} are the complete
coefficient calculus for this first inherited sector.

\section{Leading coefficient in closed form}
The coefficient of $v$ in $R_-$ comes only from $T(z^2)/2$:
\begin{equation}
 r_1^-=\frac{t_1}{\sqrt2}.
 \label{eq:rminus1}
\end{equation}
Since $z_1^-=z_0^-r_1^-$ and
$C_0'(z_0^-)=y_-/[z_0^-(1-y_-)]$,
\begin{equation}
 \boxed{
 b_1^-=\frac{y_-}{1-y_-}\frac{t_1}{\sqrt2}.}
 \label{eq:bminus1-closed}
\end{equation}
Numerically,
\begin{align}
 r_0^-={}&0.4686366279248455480872026033545827113234\ldots,\notag\\
 z_0^-={}&-0.9293757352012050888501367245727563226186\ldots,\notag\\
 y_-={}&-0.5410329414204534926205030986636417164282\ldots,\notag\\
 b_1^-={}&0.3871501110302429034772066937474883100349\ldots.
 \label{eq:minus-leading-numeric}
\end{align}
The first local coefficients are
\begin{align}
 b_1^-&= \phantom{-}0.3871501110302429034772066937\ldots,\notag\\
 b_2^-&=-0.03256491851564763414730981111\ldots,\notag\\
 b_3^-&= \phantom{-}0.01720042160312209148300754229\ldots,\notag\\
 b_4^-&= \phantom{-}0.1040770395183844385254896204\ldots,\notag\\
 b_5^-&= \phantom{-}0.1619473607727023611895041812\ldots,\notag\\
 b_6^-&=-0.06292599818799308689798687888\ldots,\notag\\
 b_7^-&=-0.2440310392314213513546115758\ldots.
 \label{eq:bminus-numeric}
\end{align}

\section{The alternating coefficient sector}
Only odd $b_{2j+1}^-$ are singular in $z$ at $\sigma_-$.  The same gamma-ratio
transfer as before gives
\begin{equation}
 a_n^{(-)}\asympseries
 \Stokes_-\sigma_-^{-n}n^{-3/2}
 \sum_{k\ge0}\frac{B_k^-}{n^k},
 \label{eq:negative-sector-general}
\end{equation}
where
\begin{equation}
 \boxed{
 B_k^-=\sum_{j=0}^{k}
 \frac{b_{2j+1}^-}{\Gamma(-j-1/2)}
 g_{k-j}\!\left(j+\frac12\right).}
 \label{eq:Bminus-master}
\end{equation}
For the direct negative-axis continuation the local Hankel multiplier is one;
we have retained $\Stokes_-$ in \cref{eq:negative-sector-general} to make the
contour dependence explicit in a larger, multi-sheet deformation.

Since $\sigma_-^{-n}=(-1)^n\rho^{-n/2}$, the numerical local sector is
\begin{align}
 a_n^{(-)}\asympseries{}&
 \Stokes_-(-1)^n\rho^{-n/2}n^{-3/2}\notag\\
 &\times\biggl(-0.1092130299563101757054885230\ldots\notag\\
 &\qquad-\frac{0.03367666220778285336719302163\ldots}{n}\notag\\
 &\qquad-\frac{0.1790009011730463163890152410\ldots}{n^2}\notag\\
 &\qquad-\frac{1.642342004953279983002352205\ldots}{n^3}\notag\\
 &\qquad-\frac{18.92664586329208396709681008\ldots}{n^4}
 -\cdots\biggr).
 \label{eq:negative-sector-numeric}
\end{align}
Relative to the leading dominant term, its first amplitude is
\begin{equation}
 \frac{B_0^-}{A_0}=-0.2482543049151649529915710045\ldots,
 \qquad
 \frac{a_n^{(-)}}{a_n^{(0)}}
 \sim -0.2482543049\ldots\,(-\sqrt\rho)^n.
 \label{eq:negative-relative}
\end{equation}
This is smaller than every power of $1/n$, and hence invisible to the complete
series \cref{eq:forward-A}; it is the first explicit genuinely transseries
correction.

\begin{remark}[Other sectors of the same or greater modulus]
The point $+\sqrt\rho$ lies behind the dominant branch point on real
continuation and belongs to sheet-dependent continuations.  Complex critical
points and overlapping root lifts may also occur farther out.  Their local
coefficients are generated by \cref{thm:local-completeness}; their global
multipliers depend on the selected contour.  Formula
\cref{eq:negative-sector-numeric} is therefore an explicit sector, not an
assertion that one term alone exhausts every continuation at modulus
$\sqrt\rho$.
\end{remark}

\chapter{Transport of exponential sectors through inversion}
\label{ch:inverse-sectors}

\section{General transseries setup}
Let the full forward interpolation be written as
\begin{equation}
 \mathcal A_{\rm full}(x)=\mathcal A_0(x)\bigl(1+\mathcal E(x)\bigr),
 \label{eq:full-forward-interpolant}
\end{equation}
where $\mathcal A_0$ contains the complete dominant inverse-power sector and
$\mathcal E$ is a well-ordered sum of exponentially small sectors,
\begin{equation}
 \mathcal E(x)=
 \sum_{\omega\in\Omega\setminus\{0\}}
 \e^{-\omega x}x^{\mu_\omega}
 \sum_{r\ge0}e_{\omega,r}x^{-r}.
 \label{eq:E-sectors}
\end{equation}
Here $\Omega$ is generated by the singular actions
$\Action_\sigma=\Log(\sigma/\rho)$.  Taking a logarithm gives
\begin{equation}
 \mathcal R(x)=\log(1+\mathcal E(x))
 =\sum_{m\ge1}\frac{(-1)^{m+1}}m\mathcal E(x)^m.
 \label{eq:R-nonpert}
\end{equation}
Products add actions, so the inverse transseries is supported on the additive
semigroup generated by $\Omega$.  Coefficients at any fixed action and inverse
power are finite Cauchy products, equivalently multivariate Bell-polynomial
expressions.

Let
\begin{equation}
 f_0(x)=\log\mathcal A_0(x),
 \qquad f_0(x_0)=\log y,
 \label{eq:f0-x0}
\end{equation}
where $x_0$ is the complete perturbative inverse of Parts~III, either in the
$W_{-1}$ form or the polynomial--logarithmic form.  Put $x=x_0+\Delta$ and
\begin{equation}
 A_{x_0}(u)=
 \begin{cases}
 \displaystyle\frac{f_0(x_0+u)-f_0(x_0)}{u},&u\ne0,\\[2mm]
 f_0'(x_0),&u=0.
 \end{cases}
 \label{eq:A-x0}
\end{equation}
Then the exact formal inversion equation is
\begin{equation}
 \Delta=\Psi_{x_0}(\Delta),
 \qquad
 \Psi_{x_0}(u)=-\frac{\mathcal R(x_0+u)}{A_{x_0}(u)}.
 \label{eq:Delta-fixed}
\end{equation}

\begin{theorem}[All-sector inverse Lagrange formula]
\label{thm:all-sector-inverse}
The beyond-all-orders correction to the inverse is
\begin{equation}
 \boxed{
 \Delta=
 \sum_{m\ge1}\frac1{m!}
 \left.\frac{\dd^{m-1}}{\dd u^{m-1}}
 \Psi_{x_0}(u)^m\right|_{u=0}.}
 \label{eq:Delta-Lagrange}
\end{equation}
When \cref{eq:R-nonpert} and $A_{x_0}^{-1}$ are expanded formally, this formula
generates every mixed action, every power of $x_0^{-1}$, and every polynomial
factor arising from differentiation.
\end{theorem}

\begin{proof}
Equation \cref{eq:Delta-fixed} is of the form required by the formal
Lagrange--B\"urmann theorem.  Introduce an auxiliary marker $t$, solve
$\Delta=t\Psi_{x_0}(\Delta)$, apply
\cref{eq:lagrange-derivative}, and set $t=1$.  Well ordering of the action
support ensures that each requested transmonomial receives only finitely many
contributions.
\end{proof}

\section{First-order transport law}
The first term of \cref{eq:Delta-Lagrange} is
\begin{equation}
 \Delta=-\frac{\mathcal R(x_0)}{f_0'(x_0)}+O(\mathcal R^2).
 \label{eq:first-sector-transport}
\end{equation}
Since $f_0'(x_0)=\lambda-p/x_0+O(x_0^{-2})$, a relative forward sector
\begin{equation}
 \mathcal E_\omega(x)
 \sim e_{\omega,0}\e^{-\omega x}x^{\mu_\omega}
 \label{eq:E-one}
\end{equation}
produces
\begin{equation}
 \Delta_\omega(y)
 \sim-\frac{e_{\omega,0}}{\lambda}
 \e^{-\omega x_0}x_0^{\mu_\omega}.
 \label{eq:Delta-one-x0}
\end{equation}
Using $y\sim C\e^{\lambda x_0}x_0^{-p}$,
\begin{equation}
 \boxed{
 \Delta_\omega(y)
 \sim-\frac{e_{\omega,0}}{\lambda}
 C^{\omega/\lambda}y^{-\omega/\lambda}
 x_0^{\mu_\omega-p\omega/\lambda}.}
 \label{eq:Delta-one-y}
\end{equation}
Complex actions therefore become complex powers of $y$ and oscillations in
$\log y$, with an additional $\log\log y$ phase modulation through $x_0$.

\section{Inverse image of the \texorpdfstring{$-\sqrt\rho$}{-sqrt(rho)} sector}
For the inherited negative sector,
\begin{equation}
 \omega_- =\Log\frac{\sigma_-}{\rho}
 =\frac\lambda2+\ii\pi
 \label{eq:omega-minus}
\end{equation}
(up to the lateral choice of $\pm\pi$), and
$e_{-,0}=B_0^-/A_0=-0.2482543049\ldots$.  A real median interpolation of the
integer alternation replaces $\e^{-\ii\pi x}$ by $\cos(\pi x)$, giving
\begin{align}
 \Delta_-(y)
 &\sim-\frac{B_0^-}{\lambda A_0}
 \e^{-\lambda x_0/2}\cos(\pi x_0)\notag\\
 &=0.2290681103840319785656467154\ldots\,
 \e^{-\lambda x_0/2}\cos(\pi x_0).
 \label{eq:inverse-minus-x0}
\end{align}
Since
$\e^{-\lambda x_0/2}\sim C^{1/2}y^{-1/2}x_0^{-3/4}$,
\begin{equation}
 \boxed{
 \Delta_-(y)
 \sim0.1519334736595309882627290391\ldots\,
 y^{-1/2}x_0^{-3/4}\cos(\pi x_0).}
 \label{eq:inverse-minus-y}
\end{equation}
The leading phase is
\begin{equation}
 \pi x_0=
 \frac{\pi}{\lambda}\log y+
 \frac{3\pi}{2\lambda}\log\log y+O(1),
 \qquad
 \frac\pi\lambda=2.898796429733978737710727372\ldots .
 \label{eq:minus-phase}
\end{equation}
Thus the first nonperturbative inverse correction has magnitude
$y^{-1/2}(\log y)^{-3/4}$ and a logarithmic oscillation whose phase contains a
slower $\log\log y$ drift.

\section{Higher corrections and mixed actions}
The next orders follow mechanically from \cref{eq:Delta-Lagrange}:
\begin{enumerate}
\item inverse powers in the $\omega_-$ sector come from $B_1^-,B_2^-,\ldots$,
      from the complete perturbative derivative $f_0'(x_0)$, and from
      differentiation of $\e^{-\omega_-u}$;
\item the two-instanton sector $2\omega_-$ is produced by the quadratic term in
      $\log(1+\mathcal E)$ and by the $m=2$ Lagrange term;
\item products with any other singularity action produce sectors
      $r\omega_-+\omega_{\sigma_1}+\cdots+\omega_{\sigma_s}$;
\item conjugate actions combine into real trigonometric polynomials in the
      corresponding logarithmic phases.
\end{enumerate}
At each fixed action and inverse order the calculation is finite.  This is the
precise sense in which \cref{eq:Delta-Lagrange} supplies general formulae for
the coefficients of the full inverse transseries.

\begin{warningbox}[title={Two distinct oscillatory phenomena}]
The exponentially small oscillation in \cref{eq:inverse-minus-y} comes from a
complex singular action and tends to zero.  The order-one sawtooth in
\cref{eq:ceiling-fourier} comes from integer rounding and never tends to zero.
They have different origins, different scales, and must not be conflated.
\end{warningbox}

\chapter{Summary of the main formulae}
\label{ch:summary}

For reference, the complete calculation may be compressed to the following
chain.

\begin{enumerate}
\item Generate $a_n$ and $b_n$ exactly from
\[
 (n-1)a_n=\sum_{j=1}^{n-1}b_ja_{n-j},
 \qquad b_j=\sum_{d\mid j}d a_d.
\]
\item Set
\[
 R(z)=\sum_{m\ge2}\frac{b_m-ma_m}{m}z^m,
 \qquad \log\rho+1+R(\rho)=0.
\]
\item Form the scaled jets
\[
 \ell_j=(-1)^{j-1}(j-1)!+\sum_{m\ge j}
 \frac{b_m-ma_m}{m}\falling mj\rho^m.
\]
\item Compute
\[
 s_j=\frac{2(-1)^{j+1}}{j!}\Bell j(\ell_1,\ldots,\ell_j),
 \qquad U(X)=\frac{\sum_{j\ge1}s_jX^j}{s_1X}.
\]
\item Use the universal coefficients
\[
 \omega_m=\frac1m[u^{m-1}]
 \left(\sum_{j\ge0}\frac{2u^j}{(j+2)j!}\right)^{-m/2}
\]
and
\[
 t_r=-\sum_{\substack{m\le r\\m\equiv r\ (2)}}
 \omega_ms_1^{m/2}[X^{(r-m)/2}]U(X)^{m/2}.
\]
\item Transfer to sequence coefficients:
\[
 A_k=\sum_{j=0}^k
 \frac{t_{2j+1}}{\Gamma(-j-1/2)}
 g_{k-j}(j+1/2),
\]
where $g_r$ is generated from Bernoulli-polynomial coefficients
\cref{eq:d-beta,eq:g-recurrence}.
\item Normalize $c_k=A_k/A_0$ and compute
\[
 \eta_n=[v^n]\log\left(1+\sum_{k\ge1}c_k\lambda^kv^k\right)
\]
by \cref{eq:eta-master}.
\item Invert with either
\[
 \nu(y)\asympseries\frac1\lambda\left(
 X_0+\sum_{r\ge1}\frac{D_r}{X_0^r}\right),
 \quad
 X_0=-\frac32W_{-1}\!\left(-\frac23
 \left(\frac{C\lambda^{3/2}}y\right)^{2/3}\right),
\]
where $D_r$ is given by \cref{eq:Dr-master}, or with the polynomial--logarithmic
recurrence \cref{eq:Pn-master}.
\item Generate all farther singular sectors recursively by
\cref{eq:regular-local-composition,eq:critical-local-composition}, transfer each
by \cref{eq:global-forward-transseries}, and invert all sectors by
\cref{eq:Delta-Lagrange}.
\end{enumerate}

\begin{resultbox}[title={Final dominant pair}]
The forward and inverse expansions, in their most useful normalized forms, are
\begin{align*}
 a_n&\asympseries
 C\alpha^n n^{-3/2}\left(1+\sum_{k\ge1}\frac{c_k}{n^k}\right),\\
 \nu(y)&\asympseries
 \frac1\lambda\left[
 Z+\frac32\log Z+
 \sum_{n\ge1}\frac{P_n(\log Z)}{Z^n}\right],
 \qquad Z=\log\frac{y}{C\lambda^{3/2}},
\end{align*}
with $C$, $\alpha$, $c_k$, and $P_n$ determined by finite coefficient formulae
above.  The first nonperturbative sector is proportional to
$(-1)^n\rho^{-n/2}n^{-3/2}$ forward and to
$y^{-1/2}(\log y)^{-3/4}$ times a logarithmic oscillation after inversion.
\end{resultbox}

\appendix

\chapter{Fully explicit finite coefficient formulae}
\label{app:finite-formulas}

The main text organizes the calculation by mathematical dependency.  This
appendix removes the remaining shorthand and displays the coefficient
operators as finite sums.  Its purpose is not to replace the more readable
formulae in the body, but to make the phrase ``general formula for the
coefficients'' completely literal.

\section{Bell-polynomial recurrences}
The partial exponential Bell polynomials satisfy
\begin{equation}
 \pBell nk(x_1,x_2,\ldots)
 =\sum_{j=1}^{n-k+1}\binom{n-1}{j-1}x_j
   \pBell{n-j}{k-1}(x_1,x_2,\ldots),
 \label{eq:bell-triangular-recurrence}
\end{equation}
with $\pBell00=1$ and zero outside $0\le k\le n$.  Thus all Bell quantities in
this article can be generated by a triangular array without enumerating
partitions.  The complete polynomials obey
\begin{equation}
 \Bell n(x_1,\ldots,x_n)
 =\sum_{j=1}^n\binom{n-1}{j-1}x_j
   \Bell{n-j}(x_1,\ldots,x_{n-j}),
 \qquad \Bell0=1.
 \label{eq:complete-bell-recurrence}
\end{equation}
Both recurrences follow by distinguishing the block containing a fixed marked
atom in the set-partition interpretation.

For a unit series $U(X)=1+\sum_{j\ge1}u_jX^j$, introduce
\begin{equation}
 \mathfrak P_N(\beta;u_1,\ldots,u_N)
 :=\frac1{N!}\sum_{k=0}^N\falling\beta k\,
 \pBell Nk(1!u_1,2!u_2,\ldots,(N-k+1)!u_{N-k+1}).
 \label{eq:Pfrak-definition}
\end{equation}
Then $\Coeff XN{U(X)^\beta}=\mathfrak P_N(\beta;u_1,\ldots,u_N)$.
In particular, every fractional power used below is an explicitly finite
polynomial in the input coefficients and in $\beta$.

\section{Universal Lambert coefficients as a partition sum}
Recall
\begin{equation}
 H(u)=\sum_{j\ge0}\frac{2u^j}{(j+2)j!}
     =1+\sum_{j\ge1}h^{\mathrm L}_ju^j,
 \qquad
 h^{\mathrm L}_j=\frac{2}{(j+2)j!}.
 \label{eq:HL-def}
\end{equation}
The superscript $\mathrm L$ distinguishes these universal Lambert coefficients
from the P\'olya-disturbance coefficients $h_j$ in \cref{eq:h-coeff}.  Combining
Lagrange inversion with \cref{eq:Pfrak-definition} gives
\begin{equation}
 \boxed{
 \omega_m=\frac{1}{m(m-1)!}
 \sum_{k=0}^{m-1}\falling{-m/2}{k}\,
 \pBell{m-1}{k}
 \bigl(1!h^{\mathrm L}_1,2!h^{\mathrm L}_2,\ldots\bigr).}
 \label{eq:omega-partition-sum}
\end{equation}
This proves at once that $\omega_m\in\mathbb Q$ for every $m$.

For convenient checking, the first eighteen values are
\begingroup
\footnotesize
\setlength{\tabcolsep}{3pt}
\begin{longtable}{rl@{\hspace{6mm}}rl}
\caption{Exact universal coefficients in $1-T=\sum_{m\ge1}\omega_mq^m$.}
\label{tab:omega-exact}\\
\toprule
$m$ & $\omega_m$ & $m$ & $\omega_m$\\
\midrule
\endfirsthead
\toprule
$m$ & $\omega_m$ & $m$ & $\omega_m$\\
\midrule
\endhead
1  & $1$ &
10 & $-5776369/1515591000$\\
2  & $-1/3$ &
11 & $169709463197/69528040243200$\\
3  & $11/72$ &
12 & $-1118511313/709296588000$\\
4  & $-43/540$ &
13 & $667874164916771/650782456676352000$\\
5  & $769/17280$ &
14 & $-500525573/744761417400$\\
6  & $-221/8505$ &
15 & $103663334225097487/234281684403486720000$\\
7  & $680863/43545600$ &
16 & $-466901817532379/1595278956070800000$\\
8  & $-1963/204120$ &
17 & $21235294185086305043/109242202556140093440000$\\
9  & $226287557/37623398400$ &
18 & $-106040742894306601/818378104464320400000$\\
\bottomrule
\end{longtable}
\endgroup

\section{One composite formula for every dominant forward coefficient}
Let
\begin{equation}
 u_j=\frac{s_{j+1}}{s_1},
 \qquad
 \mathfrak P_N(\beta)=
 \mathfrak P_N(\beta;u_1,\ldots,u_N).
 \label{eq:u-composite}
\end{equation}
Then \cref{eq:tr-master,eq:Ak-master} combine to give
\begin{align}
 A_k={}&-
 \sum_{j=0}^k\frac{g_{k-j}(j+1/2)}{\Gamma(-j-1/2)}
 \sum_{r=0}^{j}
 \omega_{2r+1}s_1^{r+1/2}
 \mathfrak P_{j-r}(r+1/2).
 \label{eq:Ak-composite}
\end{align}
No infinite formal sum occurs on the right: for fixed $k$ it contains finitely
many rational operations after the analytic jets $s_1,\ldots,s_{k+1}$ have
been supplied.  Those jets themselves are explicit convergent sums:
\begin{align}
 \ell_j
 & =(-1)^{j-1}(j-1)!+
 \sum_{m\ge j}\frac{b_m-ma_m}{m}\falling mj\rho^m,
 \label{eq:ell-composite}\\
 s_j
 & =\frac{2(-1)^{j+1}}{j!}
 \Bell j(\ell_1,\ldots,\ell_j).
 \label{eq:s-composite}
\end{align}
Consequently \cref{eq:Ak-composite,eq:ell-composite,eq:s-composite}, together
with the exact divisor recurrence \cref{eq:exact-recurrence}, are a complete
formula for $A_k$ involving no fitted constants.

\begin{remark}[Rational and analytic layers]
The numbers $\omega_m$ and the polynomials $g_r(\beta)$ are universal and
rational.  All dependence on the particular unlabelled-tree class enters
through $\rho$ and the disturbance jets $\ell_j$.  This separation is useful
both conceptually and computationally: the universal arrays may be precomputed
once and reused for any tree class satisfying $T=\zeta e^T$ with analytic
$\zeta$.
\end{remark}

\section{Completely expanded logarithmic phase coefficients}
Writing $f_j=c_j\lambda^j$, the additive phase coefficients are
\begin{equation}
 \boxed{
 \eta_n=\frac1{n!}\sum_{k=1}^n(-1)^{k-1}(k-1)!
 \pBell nk(1!f_1,2!f_2,\ldots).}
 \label{eq:eta-explicit-appendix}
\end{equation}
Equivalently, by expanding the Bell polynomial into multiplicities,
\begin{equation}
 \eta_n=
 \sum_{k=1}^{n}\frac{(-1)^{k-1}}{k}
 \sum_{\substack{m_1,m_2,\ldots\ge0\\
                  \sum_jm_j=k,\ \sum_jjm_j=n}}
 \frac{k!}{\prod_jm_j!}\prod_{j\ge1}(c_j\lambda^j)^{m_j}.
 \label{eq:eta-multiplicity}
\end{equation}
This is the ordinary logarithm formula grouped by integer partitions of $n$.

\section{A finite partition formula for the Lambert-core corrections}
Let
\begin{equation}
 \Phi_\varepsilon(u)=
 p\log(1+\varepsilon u)
 -\sum_{k\ge1}\eta_k
  \left(\frac{\varepsilon}{1+\varepsilon u}\right)^k,
 \qquad p=\frac32.
 \label{eq:Phi-eps-appendix}
\end{equation}
The coefficient $D_r$ in \cref{eq:Dr-master} may be written without derivatives
as
\begin{equation}
 \boxed{
 D_r=\sum_{m=1}^{r}\frac1m
 [\varepsilon^ru^{m-1}]\,\Phi_\varepsilon(u)^m.}
 \label{eq:Dr-bivariate}
\end{equation}
For each $r$, truncate $\Phi_\varepsilon$ at total $\varepsilon$-degree $r$;
then the coefficient in \cref{eq:Dr-bivariate} is a finite multinomial sum.
This is often the simplest exact implementation because bivariate truncated
series multiplication automatically performs all Fa\`a di Bruno bookkeeping.

\chapter{Gamma-ratio transfer to all orders}
\label{app:gamma-transfer}

\section{The exact singular-monomial coefficient}
For $\beta\notin\N_0$ and $n\ge0$,
\begin{equation}
 [z^n]\left(1-\frac z\rho\right)^\beta
 =\rho^{-n}(-1)^n\binom\beta n
 =\rho^{-n}\frac{\Gamma(n-\beta)}
 {\Gamma(-\beta)\Gamma(n+1)}.
 \label{eq:exact-transfer-appendix}
\end{equation}
If $\beta\in\N_0$, the left side vanishes for $n>\beta$.  This explains why the
integer powers $t_{2j}X^j$ in the local expansion do not contribute to the
algebraic coefficient asymptotics: they are locally analytic and have no
Hankel jump.

\section{Bernoulli-polynomial derivation}
For fixed $a,b\in\C$, Euler--Maclaurin expansion of $\log\Gamma$ gives
\begin{equation}
 \log\frac{\Gamma(n+a)}{\Gamma(n+b)n^{a-b}}
 \asympseries
 \sum_{m\ge1}\frac{(-1)^{m+1}}
 {m(m+1)}\frac{B_{m+1}(a)-B_{m+1}(b)}{n^m},
 \label{eq:log-gamma-ratio-general}
\end{equation}
where $B_m(x)$ denotes the Bernoulli polynomial.  Setting $a=-\beta$ and
$b=1$ yields the coefficients $d_m(\beta)$ of \cref{eq:d-beta}.  Exponentiating
then gives both
\begin{equation}
 g_r(\beta)=\frac1{r!}
 \Bell r(1!d_1(\beta),2!d_2(\beta),\ldots,r!d_r(\beta))
 \label{eq:g-bell-appendix}
\end{equation}
and the triangular recurrence \cref{eq:g-recurrence}.

The first five polynomials are
\begin{align}
 g_0(\beta)={}&1,\notag\\
 g_1(\beta)={}&\frac{\beta(\beta+1)}2,\notag\\
 g_2(\beta)={}&
 \frac{\beta(\beta+1)(\beta+2)(3\beta+1)}{24},\notag\\
 g_3(\beta)={}&
 \frac{\beta^2(\beta+1)^2(\beta+2)(\beta+3)}{48},\notag\\
 g_4(\beta)={}&
 \frac{\beta(\beta+1)(\beta+2)(\beta+3)(\beta+4)
 (15\beta^3+30\beta^2+5\beta-2)}{5760}.
 \label{eq:g-first-five}
\end{align}
In particular $g_r(\beta)$ is a polynomial of degree $2r$ with rational
coefficients.

\section{Remainder form of the transfer theorem}
Let $M\ge0$ be fixed.  Suppose that, in a $\Delta$-domain at $\rho$,
\begin{equation}
 T(z)=H(z)+\sum_{j=0}^{M}t_{2j+1}
 \left(1-\frac z\rho\right)^{j+1/2}
 +\bigO\!\left(\left|1-\frac z\rho\right|^{M+3/2}\right),
 \label{eq:transfer-remainder-hyp}
\end{equation}
where $H$ is analytic in a disk of radius strictly larger than $\rho$ after the
local analytic terms are patched away.  Then
\begin{equation}
 a_n=\rho^{-n}n^{-3/2}
 \left(\sum_{k=0}^{M}\frac{A_k}{n^k}
 +\bigO(n^{-M-1})\right)
 +\bigO(R^{-n})
 \label{eq:transfer-remainder-conclusion}
\end{equation}
for some $R>\rho$ determined by the chosen continuation domain.  The first
error is the algebraic truncation error at the dominant singularity; the
second records all farther singularities not crossed by the contour.  This
separation is precisely what is refined into the transseries of
\cref{eq:global-forward-transseries}.

\begin{proof}
Apply the exact identity \cref{eq:exact-transfer-appendix} to each displayed
singular monomial and truncate its gamma-ratio expansion
\cref{eq:log-gamma-ratio-general}.  A Hankel-contour estimate gives
$\rho^{-n}n^{-M-5/2}$ for the local remainder.  Cauchy's estimate on the
remaining outer contour gives $O(R^{-n})$.  This is the standard transfer
mechanism of singularity analysis \cite{FlajoletOdlyzko1990,FlajoletSedgewick2009}.
\end{proof}

\chapter{Structure and error theory of the inverse expansion}
\label{app:inverse-structure}

\section{Existence and uniqueness in the logarithmic scale}
Consider the general phase
\begin{equation}
 Z=X-p\log X+K(X^{-1}),
 \qquad
 K(v)=\sum_{j\ge1}\eta_jv^j,
 \qquad p>0.
 \label{eq:general-inverse-phase-app}
\end{equation}
The derivative is
\begin{equation}
 \frac{\dd Z}{\dd X}
 =1-\frac pX-\frac1{X^2}K'(X^{-1})=1+O(X^{-1}).
 \label{eq:phase-derivative-app}
\end{equation}
Hence the map is strictly increasing for all sufficiently large positive $X$,
and it has a unique inverse germ at $+
ealsymbol$.

\begin{theorem}[Formal uniqueness]
There is a unique logarithmic transseries of the form
\begin{equation}
 X=Z+p\log Z+\sum_{n\ge1}P_n(\log Z)Z^{-n}
 \label{eq:formal-unique-inverse}
\end{equation}
that solves \cref{eq:general-inverse-phase-app} coefficientwise.  Each $P_n$ is
a polynomial with coefficients in
$\mathbb Q[p,\eta_1,\ldots,\eta_n]$.
\end{theorem}

\begin{proof}
Substitute \cref{eq:formal-unique-inverse}, expand
$\log X=\log Z+\log(1+(X-Z)/Z)$ and expand $K(X^{-1})$ as a unit-series
composition.  At order $Z^{-n}$ the unknown $P_n$ occurs linearly with
coefficient one, while all other terms involve only $P_0,\ldots,P_{n-1}$.
Thus the coefficients are determined triangularly.  The operations used are
addition, multiplication, formal logarithm, and unit-series powers, so they
preserve the stated coefficient ring.
\end{proof}

\section{Degree, leading term, and constant term}
The triangular recurrence contains additional structure.

\begin{proposition}[Shape of the inverse polynomials]
For $n\ge1$,
\begin{equation}
 \deg P_n=n,
 \qquad
 [L^n]P_n(L)=\frac{(-1)^{n+1}p^{n+1}}{n},
 \qquad
 P_n(0)=D_n.
 \label{eq:Pn-shape}
\end{equation}
Here $D_n$ is the Lambert-core coefficient from \cref{eq:Dr-master}.
\end{proposition}

\begin{proof}
The highest power of $L$ is independent of $K$: every occurrence of an
$\eta_j$ consumes at least one inverse power without increasing the maximum
logarithmic degree.  Thus the leading term is the same as for
$Z=X-p\log X$.  Expanding its exact inverse
$X=-pW_{-1}(-e^{-Z/p}/p)$, or inducting directly in
\cref{eq:Pn-master}, gives $(-1)^{n+1}p^{n+1}L^n/n$.

For the constant term, set $L=0$ in the direct recurrence.  The resulting
near-identity reversion is exactly the correction equation about the Lambert
core with $X_0^{-1}$ as small parameter.  Formal uniqueness therefore identifies
its order-$n$ coefficient with $D_n$.
\end{proof}

\section{Transport of a controlled forward remainder}
Suppose a positive smooth interpolant has logarithm
\begin{equation}
 f(x)=\lambda x-p\log x+\log C+
 \sum_{j=1}^{M}\kappa_jx^{-j}+O(x^{-M-1}),
 \qquad \lambda>0.
 \label{eq:forward-log-remainder}
\end{equation}
Let $x_M(y)$ be the inverse obtained by matching the displayed expansion through
order $M$.  Since $f'(x)=\lambda+O(x^{-1})$, the mean-value theorem gives
\begin{equation}
 x_M(y)-f^{-1}(\log y)
 =O\!\left((\log y)^{-M-1}\right).
 \label{eq:inverse-remainder-transport}
\end{equation}
More precisely, the leading residual is divided by
$f'(x_M)=\lambda-p/x_M+O(x_M^{-2})$.  This explains why forward relative errors
and inverse absolute errors in \cref{tab:inverse-errors} have nearly the same
asymptotic order, apart from the factor $1/\lambda$.

\section{Staircase recovery and the best possible deterministic error}
Let $\nu$ be any continuous asymptotic inverse agreeing with a monotone
interpolation through the sequence values.  Then
\begin{equation}
 N_*(y)=\lceil \nu(y)+\epsilon(y)\rceil,
 \qquad \epsilon(y)\to0,
 \label{eq:staircase-general}
\end{equation}
provided $y$ stays away from the finitely many initial ties.  No continuous
asymptotic series can determine the integer answer with an error tending to
zero uniformly in $y$: arbitrarily small changes of $y$ across $a_n$ change
$N_*(y)$ by one.  The Fourier sawtooth in \cref{eq:ceiling-fourier} is therefore
not an optional refinement but the canonical order-one transseries layer of
the generalized inverse.

\chapter{Numerical audit and reusable constants}
\label{app:numerical-audit}

\section{Dominant singular constants}
The constants used in the tables are
\begin{align}
 \rho={}&0.3383218568992076951961126257170170531837746075329677955723037763
 \ldots,\notag\\
 \alpha={}&2.955765285651994974714817524123194588375492304663596595350472479
 \ldots,\notag\\
 \lambda={}&1.083757597244624660482711980896174279049662124202672234914427065
 \ldots,\notag\\
 C=A_0={}&0.43992401257102530404090339143454476479808540794012\ldots,\notag\\
 C\lambda^{3/2}={}&0.49633614165959974747273694642522491348\ldots .
 \label{eq:numerical-audit-constants}
\end{align}
They satisfy the independent checks
\begin{equation}
 \lambda=-\log\rho=1+R(\rho),
 \qquad
 C^2=\frac{1+\rho R'(\rho)}{2\pi}.
 \label{eq:constant-checks}
\end{equation}

\section{Local Puiseux coefficients}
For $X=1-z/\rho$,
$T(z)=1+\sum_{r\ge1}t_rX^{r/2}$ begins as follows.
\begin{longtable}{r@{\qquad}r@{\qquad}r@{\qquad}r}
\caption{Numerical dominant Puiseux coefficients.}
\label{tab:t-audit}\\
\toprule
$r$ & $t_r$ & $r$ & $t_r$\\
\midrule
\endfirsthead
\toprule
$r$ & $t_r$ & $r$ & $t_r$\\
\midrule
\endhead
1  & $-1.559490020374640885542207396$ &
10 & $ 0.2059848312779074187361657348$\\
2  & $ 0.8106697078826992814381417462$ &
11 & $-0.05272470849819056231533821627$\\
3  & $-0.2854870216128456053038513708$ &
12 & $ 0.01702656875495366944604539502$\\
4  & $ 0.1653723657120838934765135390$ &
13 & $-0.1523706243663253962860946421$\\
5  & $-0.3424599704021542010422002117$ &
14 & $ 0.1737028832998504628939446698$\\
6  & $ 0.3174072259465285635680708874$ &
15 & $-0.01447370373952704485673111250$\\
7  & $-0.1077788002916310091822310895$ &
16 & $-0.02189951761121556223267283747$\\
8  & $ 0.06138495705583510468873149925$ &
17 & $-0.1445471935709097054769461137$\\
9  & $-0.1952123835975564636127698720$ &
18 & $ 0.1760771088850177790623322866$\\
\bottomrule
\end{longtable}
The alternating, nonmonotone pattern is expected: these are local analytic
coefficients, not the final coefficient-asymptotic constants, and even-indexed
terms do not directly survive singularity transfer.

\section{Inverse phase and Lambert-core coefficients}
The first additive phase coefficients $\eta_j$ and correction coefficients
$D_j$ are
\begin{longtable}{r@{\quad}r@{\qquad}r}
\caption{Coefficients in $K(v)=\sum\eta_jv^j$ and
$\delta(X_0)=\sum D_jX_0^{-j}$.}
\label{tab:eta-D-audit}\\
\toprule
$j$ & $\eta_j$ & $D_j$\\
\midrule
\endfirsthead
\toprule
$j$ & $\eta_j$ & $D_j$\\
\midrule
\endhead
1 & $0.1088130343442745145167554784$ & $-0.1088130343442745145167554784$\\
2 & $0.5859660997717814776889432009$ & $-0.7491856512881932494640764185$\\
3 & $2.446558573456779649881479860$  & $-3.582177326832123\ldots$\\
4 & $14.06753235112498242303023266$  & $-19.658721211387\ldots$\\
5 & $106.8546611534523691271119951$  & $-138.532747851856\ldots$\\
6 & $1022.176467223247872767588518$  & $-1248.97932637169\ldots$\\
7 & $11859.388275\ldots$              & $-13905.4088447481\ldots$\\
8 & $162087.382638\ldots$             & $-184719.191126508\ldots$\\
\bottomrule
\end{longtable}
The growth visible here is one reason to retain the exact $W_{-1}$ core: it
sums the entire logarithmic skeleton before the rapidly growing correction
series is truncated.

\section{Exact integer benchmarks}
The divisor recurrence gives, without floating-point arithmetic,
\begin{center}
\begin{tabular}{r@{\qquad}r}
\toprule
$n$ & $a_n$\\
\midrule
1   & $1$\\
10  & $719$\\
20  & $12\,826\,228$\\
30  & $354\,426\,847\,597$\\
40  & $11\,703\,780\,079\,612\,453$\\
50  & $425\,976\,989\,835\,141\,038\,353$\\
75  & $135\,410\,903\,503\,191\,503\,503\,384\,705\,970\,501$\\
100 & $51\,384\,328\,351\,659\,326\,880\,337\,136\,395\,054\,298\,255\,277\,970$\\
\bottomrule
\end{tabular}
\end{center}
These values provide direct regression tests for any implementation of
\cref{eq:exact-recurrence} and for the error tables in the main text.

\chapter{Notation and scale map}
\label{app:notation}

Several variables deliberately carry similar letters because each normalizes a
different stage.  The following table collects them in one place.
\begin{longtable}{p{0.16\textwidth}p{0.72\textwidth}}
\caption{Principal symbols and their roles.}
\label{tab:notation}\\
\toprule
Symbol & Meaning\\
\midrule
\endfirsthead
\toprule
Symbol & Meaning\\
\midrule
\endhead
$a_n$ & Number of unlabelled non-plane rooted trees on $n$ vertices; A000081.\\
$T(z)$ & Ordinary generating function $\sum_{n\ge1}a_nz^n$.\\
$b_m$ & Divisor transform $\sum_{d\mid m}d a_d$.\\
$R(z)$ & Analytic disturbance $\sum_{k\ge2}T(z^k)/k$.\\
$\zeta(z)$ & Cayley argument $z\exp R(z)$, so $T=-W_0(-\zeta)$.\\
$\rho$ & Positive dominant singularity; $T(\rho)=1$.\\
$\alpha,\lambda$ & $\alpha=\rho^{-1}$ and $\lambda=\log\alpha$.\\
$X=1-z/\rho$ & Local forward singular coordinate.\\
$q$ & Universal Lambert coordinate $\sqrt{2(1-e\zeta(z))}$.\\
$\omega_m$ & Rational inverse coefficients in $1-T=\sum\omega_mq^m$.\\
$t_r$ & Puiseux coefficients in $T=1+\sum t_rX^{r/2}$.\\
$A_k,c_k$ & Absolute and relative dominant coefficient-asymptotic constants.\\
$x$ & Continuous large index in the smooth interpolation $\mathcal A(x)$.\\
$X=\lambda x$ & Dimensionless inverse index; context distinguishes it from the local forward $X$.\\
$Z$ & Normalized logarithm $\log(y/(C\lambda^p))$.\\
$L$ & Slow logarithm $\log Z$.\\
$K(v)$ & Additive correction phase $\log F(\lambda v)$.\\
$X_0$ & Exact Lambert core $-pW_{-1}(-e^{-Z/p}/p)$.\\
$D_r$ & Ordinary inverse-power corrections about $X_0$.\\
$P_r(L)$ & Polynomial--logarithmic coefficients in the direct inverse expansion.\\
$\sigma,\Action_\sigma$ & A farther singularity and its action $\Log(\sigma/\rho)$.\\
$\Stokes_\sigma$ & Contour intersection/Stokes multiplier for a chosen continuation sheet.\\
\bottomrule
\end{longtable}

\backmatter

\chapter*{Bibliographic note}
\addcontentsline{toc}{chapter}{Bibliographic note}
P\'olya's cycle-index method and Otter's analysis supply the classical rooted-tree
foundation.  The Cayley-disturbance viewpoint and full dominant expansion are
closest in spirit to Genitrini's treatment, while the explicit coefficient
organization here follows the Bell-polynomial and reversion calculus developed
in the linked ProveIt drafts.  The all-order inverse, the distinction between
smooth and staircase inverses, the recursively contour-resolved farther
singularity sectors, and the explicit transport of those sectors through the
inverse are developed here as a unified extension of those ingredients.

\begin{thebibliography}{99}

\bibitem{BellBurrisYeats2006}
J.~P. Bell, S.~N. Burris, and K.~A. Yeats,
``Counting rooted trees: the universal law
$t(n)\sim C\rho^{-n}n^{-3/2}$,''
\emph{Electronic Journal of Combinatorics} \textbf{13} (2006), Research Paper
R63.

\bibitem{BergeronLabelleLeroux1998}
F.~Bergeron, G.~Labelle, and P.~Leroux,
\emph{Combinatorial Species and Tree-Like Structures},
Encyclopedia of Mathematics and its Applications, vol.~67,
Cambridge University Press, 1998.

\bibitem{Butcher1963}
J.~C. Butcher,
``Coefficients for the study of Runge--Kutta integration processes,''
\emph{Journal of the Australian Mathematical Society} \textbf{3} (1963),
185--201.

\bibitem{Butcher1972}
J.~C. Butcher,
``An algebraic theory of integration methods,''
\emph{Mathematics of Computation} \textbf{26} (1972), 79--106.

\bibitem{Butcher2021}
J.~C. Butcher,
\emph{B-Series: Algebraic Analysis of Numerical Methods},
Springer Series in Computational Mathematics, vol.~55, Springer, 2021.

\bibitem{Comtet1974}
L.~Comtet,
\emph{Advanced Combinatorics: The Art of Finite and Infinite Expansions},
D. Reidel, 1974.

\bibitem{Corless1996}
R.~M. Corless, G.~H. Gonnet, D.~E.~G. Hare, D.~J. Jeffrey, and D.~E. Knuth,
``On the Lambert $W$ function,''
\emph{Advances in Computational Mathematics} \textbf{5} (1996), 329--359.

\bibitem{DLMF}
NIST Digital Library of Mathematical Functions,
Chapter~5, ``Gamma Function,'' especially the asymptotic-expansion sections,
\url{https://dlmf.nist.gov/5}.

\bibitem{Drmota2009}
M.~Drmota,
\emph{Random Trees: An Interplay between Combinatorics and Probability},
Springer, 2009.

\bibitem{DrmotaGittenberger2010}
M.~Drmota and B.~Gittenberger,
``The shape of unlabeled rooted random trees,''
\emph{European Journal of Combinatorics} \textbf{31}(8) (2010), 2028--2063.

\bibitem{FlajoletOdlyzko1990}
P.~Flajolet and A.~Odlyzko,
``Singularity analysis of generating functions,''
\emph{SIAM Journal on Discrete Mathematics} \textbf{3}(2) (1990), 216--240.

\bibitem{FlajoletSedgewick2009}
P.~Flajolet and R.~Sedgewick,
\emph{Analytic Combinatorics}, Cambridge University Press, 2009.

\bibitem{Genitrini2016}
A.~Genitrini,
``Full asymptotic expansion for P\'olya structures,''
Proceedings of the 27th International Meeting on Probabilistic,
Combinatorial and Asymptotic Methods for the Analysis of Algorithms,
Krakow, 2016; arXiv:1605.00837.

\bibitem{McLachlan2017}
R.~I. McLachlan, K.~Modin, H.~Munthe-Kaas, and O.~Verdier,
``Butcher series: a story of rooted trees and numerical methods for evolution
equations,'' \emph{Asia Pacific Mathematics Newsletter} \textbf{7}(1) (2017),
1--11.

\bibitem{OEISA000081}
OEIS Foundation Inc.,
``A000081: number of unlabeled rooted trees with $n$ nodes,''
\url{https://oeis.org/A000081}, accessed September 2026.

\bibitem{Olver1997}
F.~W.~J. Olver,
\emph{Asymptotics and Special Functions}, A K Peters, 1997 reprint.

\bibitem{Otter1948}
R.~Otter,
``The number of trees,''
\emph{Annals of Mathematics} \textbf{49}(3) (1948), 583--599.

\bibitem{Polya1937}
G.~P\'olya,
``Kombinatorische Anzahlbestimmungen f\"ur Gruppen, Graphen und chemische
Verbindungen,'' \emph{Acta Mathematica} \textbf{68} (1937), 145--254.

\bibitem{ReshetnikovCCC}
V.~Reshetnikov,
\emph{Combinatorial Coefficient Calculus},
ProveIt semi-formalized research-frontier draft, 2026,
\url{https://github.com/VladimirReshetnikov/ProveIt/tree/main/Analysis/FabiusFunction/docs/semi-formalized-research-frontiers/drafts/combinatorial-coefficient-calculus/Combinatorial_Coefficient_Calculus}.

\bibitem{ReshetnikovSeries}
V.~Reshetnikov,
\emph{Series and Transseries},
ProveIt semi-formalized research-frontier draft collection, 2026,
\url{https://github.com/VladimirReshetnikov/ProveIt/tree/main/Analysis/FabiusFunction/docs/semi-formalized-research-frontiers/drafts/series-and-transseries}.

\bibitem{TricomiErdelyi1951}
F.~G. Tricomi and A.~Erd\'elyi,
``The asymptotic expansion of a ratio of gamma functions,''
\emph{Pacific Journal of Mathematics} \textbf{1}(1) (1951), 133--142.

\end{thebibliography}

\end{document}
